% !TEX program = xelatex
\documentclass[a4paper,12pt]{extreport}

% ============================================================
% FONT, LANGUAGE, AND PAGE FORMAT
% ============================================================
\usepackage{fontspec}
\defaultfontfeatures{Ligatures=TeX}
\IfFontExistsTF{Times New Roman}{
  \setmainfont{Times New Roman}
}{
  \setmainfont{Tinos}
}
\setmonofont{DejaVu Sans Mono}
\usepackage[vietnamese]{babel}
\makeatletter
\renewcommand\normalsize{%
  \@setfontsize\normalsize{13pt}{15.6pt}%
  \abovedisplayskip 10pt plus 2pt minus 5pt%
  \abovedisplayshortskip 0pt plus 3pt%
  \belowdisplayshortskip 6pt plus 3pt minus 3pt%
  \belowdisplayskip \abovedisplayskip%
  \let\@listi\@listI}
\makeatother
\normalsize
\usepackage[
  paperwidth=210mm,
  paperheight=297mm,
  left=3.5cm,
  right=2cm,
  top=2cm,
  bottom=2cm,
  headsep=0.5cm,
  footskip=1cm
]{geometry}
\usepackage{setspace}
\setstretch{1.2}
\setlength{\parindent}{1.5cm}
\setlength{\parskip}{0pt}
\usepackage{indentfirst}
\usepackage{microtype}
\emergencystretch=2em

% ============================================================
% HEADINGS AND PAGINATION
% ============================================================
\usepackage{titlesec}
\titleformat{\chapter}[display]
  {\normalfont\fontsize{16}{19}\selectfont\bfseries\centering}
  {CHƯƠNG \thechapter}
  {12pt}
  {\MakeUppercase}
\titlespacing*{\chapter}{0pt}{0pt}{18pt}

\titleformat{\section}
  {\normalfont\fontsize{14}{17}\selectfont\bfseries}
  {\thesection}{0.5em}{}
\titlespacing*{\section}{0pt}{6pt}{0pt}

\titleformat{\subsection}
  {\normalfont\fontsize{13}{16}\selectfont\bfseries}
  {\thesubsection}{0.5em}{}
\titlespacing*{\subsection}{0pt}{6pt}{0pt}

\titleformat{\subsubsection}
  {\normalfont\fontsize{13}{16}\selectfont\bfseries\itshape}
  {\thesubsubsection}{0.5em}{}
\titlespacing*{\subsubsection}{0pt}{6pt}{0pt}

\usepackage{needspace}
\usepackage{fancyhdr}
\pagestyle{fancy}
\fancyhf{}
\fancyfoot[C]{\fontsize{13}{15}\selectfont\thepage}
\renewcommand{\headrulewidth}{0pt}
\renewcommand{\footrulewidth}{0pt}
\setlength{\headheight}{14.5pt}

% ============================================================
% TABLES, FIGURES, MATH, LISTINGS
% ============================================================
\usepackage{amsmath,amssymb}
\usepackage{array,multirow,booktabs,tabularx,longtable}
\usepackage{makecell}
\usepackage{graphicx}
\graphicspath{{images/}}
\usepackage{float}
\usepackage{pdflscape}
\usepackage{caption}
\captionsetup[table]{
  labelsep=colon,
  font={normalsize},
  skip=3pt,
  position=above,
  justification=raggedright,
  singlelinecheck=false
}
\captionsetup[figure]{
  labelsep=colon,
  font={normalsize},
  skip=3pt,
  position=bottom,
  justification=centering
}
\usepackage{tikz}
\usetikzlibrary{arrows.meta,positioning,shapes.geometric,fit,calc}
\usepackage{xcolor}
\definecolor{DraftYellow}{RGB}{255,248,204}
\definecolor{DraftBorder}{RGB}{150,120,0}
\definecolor{SoftBlue}{RGB}{235,243,252}
\definecolor{SoftGray}{RGB}{245,245,245}
\definecolor{DarkBlue}{RGB}{31,78,121}

\usepackage{listings}
\lstdefinelanguage{json}{
  basicstyle=\ttfamily\fontsize{10}{12}\selectfont,
  numbers=left,
  numberstyle=\tiny,
  stepnumber=1,
  numbersep=7pt,
  showstringspaces=false,
  breaklines=true,
  frame=single,
  literate=
   *{0}{{{\color{blue}0}}}{1}
    {1}{{{\color{blue}1}}}{1}
    {2}{{{\color{blue}2}}}{1}
    {3}{{{\color{blue}3}}}{1}
    {4}{{{\color{blue}4}}}{1}
    {5}{{{\color{blue}5}}}{1}
    {6}{{{\color{blue}6}}}{1}
    {7}{{{\color{blue}7}}}{1}
    {8}{{{\color{blue}8}}}{1}
    {9}{{{\color{blue}9}}}{1}
    {:}{{{\color{black}:}}}{1}
    {,}{{{\color{black},}}}{1}
    {\{}{{{\color{black}\{}}}{1}
    {\}}{{{\color{black}\}}}}{1}
    {[}{{{\color{black}[}}}{1}
    {]}{{{\color{black}]}}}{1},
}
\lstset{
  basicstyle=\ttfamily\fontsize{10}{12}\selectfont,
  breaklines=true,
  frame=single,
  numbers=left,
  numberstyle=\tiny,
  numbersep=7pt,
  showstringspaces=false,
  tabsize=2,
  xleftmargin=0.6cm,
  framexleftmargin=0.4cm
}

% ============================================================
% TOC, LINKS, BIBLIOGRAPHY
% ============================================================
\usepackage{tocloft}
\renewcommand{\cftchapfont}{\bfseries}
\renewcommand{\cftsecfont}{\normalfont}
\renewcommand{\cftsubsecfont}{\normalfont}
\renewcommand{\cftchapdotsep}{\cftdotsep}
\setcounter{tocdepth}{3}
\setcounter{secnumdepth}{3}

\usepackage{xurl}
\usepackage[hidelinks,unicode]{hyperref}
\usepackage{bookmark}
\hypersetup{
  pdftitle={A Structure-Aware and Temporal RAG System for Vietnamese Traffic Law Question Answering with Verifiable Citations},
  pdfauthor={Quach Truong Phuc},
  pdfsubject={Khoa luan tot nghiep},
  pdfkeywords={RAG, Vietnamese law, temporal retrieval, citation verification, structure-aware, evidence completeness}
}

\usepackage{csquotes}
\DeclareQuoteStyle{vietnamese}{“}{”}{‘}{’}
\usepackage[
  backend=biber,
  style=numeric,
  sorting=none,
  giveninits=true,
  maxbibnames=20,
  doi=true,
  url=true
]{biblatex}
\addbibresource{references.bib}
\DeclareLanguageMapping{vietnamese}{english}
\AtBeginBibliography{\small}

% ============================================================
% DOCUMENT MACROS
% ============================================================
\newif\ifdraftmode
\draftmodetrue

\newcommand{\StudentName}{Quách Trương Phúc}
\newcommand{\StudentId}{[Mã số sinh viên]}
\newcommand{\AdvisorName}{[Tên giảng viên hướng dẫn]}
\newcommand{\ThesisMonth}{9}
\newcommand{\ThesisYear}{2026}
\newcommand{\ThesisTitleVN}{XÂY DỰNG HỆ THỐNG RAG NHẬN BIẾT CẤU TRÚC VÀ THỜI GIAN HIỆU LỰC ĐỂ HỖ TRỢ TRA CỨU PHÁP LUẬT GIAO THÔNG VIỆT NAM VỚI TRÍCH DẪN CÓ THỂ KIỂM CHỨNG}
\newcommand{\ThesisTitleEN}{A Structure-Aware and Temporal RAG System for Vietnamese Traffic Law Question Answering with Verifiable Citations}

\newcommand{\draftnotice}{%
  \ifdraftmode
  \begin{center}
  \fcolorbox{DraftBorder}{DraftYellow}{%
    \parbox{0.90\textwidth}{%
      \centering\bfseries
      BẢN THẢO KỸ THUẬT - CHƯA CÓ KẾT QUẢ THỰC NGHIỆM CUỐI.\\
      Các bảng kết quả trong Chương 6 phải được thay bằng số liệu từ evaluation run đã đóng băng trước khi nộp.
    }
  }
  \end{center}
  \fi
}

\newcommand{\pendingresult}[1]{%
  \ifdraftmode
  \begin{center}
  \fcolorbox{DraftBorder}{DraftYellow}{%
    \parbox{0.90\textwidth}{\textbf{Nội dung chờ cập nhật:} #1}
  }
  \end{center}
  \else
  #1
  \fi
}

\newcommand{\placeholderfigure}[2][0.80\textwidth]{%
  \begin{figure}[H]
    \centering
    \fbox{\parbox[c][5cm][c]{#1}{\centering\textit{Vị trí chèn hình sau khi hoàn thiện giao diện hoặc xuất diagram chính thức}}}
    \caption{#2}
  \end{figure}
}

\newcommand{\code}[1]{\texttt{#1}}
\newcolumntype{Y}{>{\raggedright\arraybackslash}X}
\newcommand{\apipath}[1]{\path{#1}}
\renewcommand{\arraystretch}{1.18}

% ============================================================
\begin{document}

% ============================================================
% MAIN COVER
% ============================================================
\thispagestyle{empty}
\begin{center}
{\fontsize{16}{19}\selectfont\bfseries
BỘ GIÁO DỤC VÀ ĐÀO TẠO\\
TRƯỜNG ĐẠI HỌC TÂY ĐÔ\\
KHOA KỸ THUẬT - CÔNG NGHỆ\\[6pt]
}

\IfFileExists{images/logo-tdu.png}{%
  \includegraphics[width=2.6cm]{images/logo-tdu.png}\\[4pt]
}{%
  \vspace{1.3cm}
}

{\fontsize{13}{16}\selectfont\bfseries
Trí tuệ - Sáng tạo - Năng động - Đổi mới
}

\vspace{1.8cm}

{\fontsize{16}{19}\selectfont\bfseries KHÓA LUẬN TỐT NGHIỆP}

\vspace{1.8cm}

{\fontsize{18}{22}\selectfont\bfseries \ThesisTitleVN}

\vspace{1.8cm}

\begin{tabular}{>{\bfseries}rl}
Chuyên ngành: & Công nghệ thông tin\\
Ngành: & Công nghệ thông tin\\
\end{tabular}

\vspace{1.6cm}

\begin{tabular}{>{\bfseries}rl}
Sinh viên thực hiện: & \StudentName\\
Mã số sinh viên: & \StudentId\\
Người hướng dẫn: & \AdvisorName\\
\end{tabular}

\vfill
{\fontsize{14}{17}\selectfont Cần Thơ, tháng \ThesisMonth\ năm \ThesisYear}
\end{center}
\clearpage

% ============================================================
% INNER COVER
% ============================================================
\thispagestyle{empty}
\begin{center}
{\fontsize{16}{19}\selectfont\bfseries
BỘ GIÁO DỤC VÀ ĐÀO TẠO\\
TRƯỜNG ĐẠI HỌC TÂY ĐÔ\\
KHOA KỸ THUẬT - CÔNG NGHỆ
}

\vspace{2.2cm}

{\fontsize{16}{19}\selectfont\bfseries KHÓA LUẬN TỐT NGHIỆP}

\vspace{1.8cm}

{\fontsize{18}{22}\selectfont\bfseries \ThesisTitleVN}

\vspace{1cm}

{\fontsize{13}{16}\selectfont\itshape \ThesisTitleEN}

\vspace{1.8cm}

\begin{tabular}{>{\bfseries}rl}
Chuyên ngành: & Công nghệ thông tin\\
Ngành: & Công nghệ thông tin\\
Sinh viên thực hiện: & \StudentName\\
Mã số sinh viên: & \StudentId\\
Người hướng dẫn: & \AdvisorName\\
\end{tabular}

\vfill
{\fontsize{14}{17}\selectfont Cần Thơ, tháng \ThesisMonth\ năm \ThesisYear}
\end{center}
\clearpage

% ============================================================
% FRONT MATTER
% ============================================================
\pagenumbering{roman}
\setcounter{page}{1}

\chapter*{Lời Cảm Ơn}
\addcontentsline{toc}{chapter}{Lời cảm ơn}

Tác giả xin gửi lời cảm ơn chân thành đến \AdvisorName\ đã dành thời gian hướng dẫn, góp ý và hỗ trợ trong quá trình xác định bài toán, lựa chọn phương pháp và xây dựng hệ thống. Các nhận xét chuyên môn đã giúp đề tài được điều chỉnh từ một chatbot hỏi đáp tài liệu tổng quát thành một nghiên cứu có phạm vi rõ ràng về truy xuất pháp luật giao thông theo cấu trúc, quan hệ tham chiếu và thời gian hiệu lực.

Tác giả trân trọng cảm ơn quý thầy cô Khoa Kỹ thuật - Công nghệ, Trường Đại học Tây Đô đã cung cấp nền tảng kiến thức và tạo điều kiện để khóa luận được thực hiện. Tác giả cũng cảm ơn gia đình, bạn bè và những người đã hỗ trợ kiểm tra tài liệu, đối chiếu nội dung văn bản pháp luật, trao đổi tình huống sử dụng và động viên trong thời gian triển khai đề tài.

Khóa luận còn có thể tồn tại thiếu sót do giới hạn về thời gian, quy mô corpus và kinh nghiệm nghiên cứu. Tác giả mong nhận được ý kiến góp ý để hệ thống và phương pháp đánh giá tiếp tục được hoàn thiện.
\clearpage

\chapter*{Lời Cam Đoan}
\addcontentsline{toc}{chapter}{Lời cam đoan}

Tôi xin cam đoan rằng khóa luận này là do chính tôi thực hiện, không sao chép dưới bất kỳ hình thức nào hay thuê hoặc nhờ người khác thực hiện.

Dữ liệu và kết quả phân tích trong khóa luận bảo đảm tính chính xác, khách quan và trung thực, không có bất kỳ sự ngụy tạo hoặc điều chỉnh kết quả nghiên cứu theo chủ quan của tác giả. Những nội dung, phương pháp và tài liệu được kế thừa từ công trình khác đều được trích dẫn trong danh mục tài liệu tham khảo.

\vspace{36pt}
\begin{flushright}
\begin{tabular}{c}
Cần Thơ, ngày \ldots\ tháng \ldots\ năm 2026\\[24pt]
Tác giả\\[24pt]
(Ký tên và ghi rõ họ và tên)\\[18pt]
\StudentName
\end{tabular}
\end{flushright}
\clearpage

\chapter*{Tóm Tắt}
\addcontentsline{toc}{chapter}{Tóm tắt}

Khóa luận nghiên cứu và xây dựng một hệ thống Retrieval-Augmented Generation (RAG) hỗ trợ tra cứu pháp luật giao thông đường bộ Việt Nam. Khác với mô hình hỏi đáp tài liệu thông thường, hệ thống được thiết kế để nhận biết cấu trúc Chương, Mục, Điều, Khoản và Điểm kèm nhãn Điểm tiếng Việt (a) b) c) d) đ) e)); mô hình hóa quan hệ tham chiếu chéo giữa các quy định; quản lý khoảng thời gian hiệu lực của từng provision; đồng thời kiểm chứng trích dẫn theo nhiều tầng trước khi trả kết quả.

Pipeline ingestion sử dụng Parser Router định tuyến tài liệu qua Docling hoặc MinerU theo đặc tính tài liệu và quality gate, sau đó chuẩn hóa sang Canonical Document IR do dự án sở hữu. Legal Structure Extractor trích xuất phân cấp pháp luật Việt Nam và tạo định danh provision ổn định; Legal Reference Resolver và Temporal and Amendment Resolver xác định quan hệ tham chiếu và khoảng hiệu lực. PostgreSQL là nguồn chân lý dữ liệu pháp lý; Qdrant là index dẫn xuất lưu dense vector, sparse BM25 và payload phục vụ truy xuất lai. Ingestion chạy nền qua Redis và Dramatiq; PDF nguồn và artifact được lưu trong MinIO.

Pipeline hỏi đáp sử dụng LangGraph như một workflow có đường đi xác định trước, gồm query understanding, temporal resolution, query expansion, retrieval đa kênh (exact lookup, dense, sparse) với RRF fusion, reranking, mở rộng ngữ cảnh pháp lý theo quan hệ và Evidence Completeness Gate. Gemini 3.5 Flash sinh câu trả lời có cấu trúc cấp claim, trong đó mỗi claim chỉ được tham chiếu \code{provision\_id} thuộc context. Bộ kiểm chứng sáu tầng (schema, citation ID, temporal, numeric grounding, claim support, evidence completeness) quyết định trả verified answer hay abstention với lý do chuẩn. Khi bằng chứng không đủ, hệ thống từ chối trả lời thay vì tìm kiếm web tự do.

\begin{sloppypar}
Nghiên cứu đề xuất gold set 200 câu đã review, chia 40 development, 40 validation và 120 final test, cùng bốn suite thí nghiệm A-D: parser benchmark (P1-P3), embedding benchmark (E1-E3), retrieval ablation (R1-R10) và generation/verification ablation (G1-G7). Các metric chính gồm Recall@k, MRR@10, nDCG@10, Citation Precision/Recall/F1, Temporal Validity Accuracy, Numeric Grounding Accuracy, Evidence Completeness và Abstention F1. Langfuse phục vụ observability ngoài đường tới hạn; RAGFlow chỉ được dùng làm baseline so sánh bên ngoài.
\end{sloppypar}

\draftnotice

\noindent\textbf{Từ khóa:} Retrieval-Augmented Generation, pháp luật giao thông Việt Nam, truy xuất theo thời gian, hybrid retrieval, kiểm chứng trích dẫn, evidence completeness, Parser Router, Canonical Document IR.
\clearpage

\chapter*{Abstract}
\addcontentsline{toc}{chapter}{Abstract}

This thesis investigates and develops a Retrieval-Augmented Generation system for Vietnamese road traffic law question answering. Unlike a conventional document chatbot, the proposed system is designed to preserve the legal hierarchy of chapters, sections, articles, clauses, and points including Vietnamese point labels (a) b) c) d) đ) e)); to model cross-references between provisions; to represent the effective period of each provision; and to verify citations through multiple layers before returning an answer.

The offline ingestion pipeline uses a Parser Router that dispatches documents to Docling or MinerU based on document characteristics and quality gates, then normalizes the output into a project-owned Canonical Document IR. A dedicated Legal Structure Extractor recognizes the Vietnamese legal hierarchy and generates stable provision identifiers; a Legal Reference Resolver and a Temporal and Amendment Resolver derive cross-reference relations and effective intervals. PostgreSQL is the legal source of truth, while Qdrant is a rebuildable retrieval index holding dense vectors, sparse BM25, and payload metadata. Ingestion runs in the background with Redis and Dramatiq; source PDFs and artifacts are stored in MinIO.

The online query pipeline is a controlled LangGraph workflow with predefined branches: query understanding, temporal resolution, query expansion, multi-channel retrieval (exact lookup, dense, sparse) fused by Reciprocal Rank Fusion, reranking, relation-aware legal context expansion, and an Evidence Completeness Gate. Gemini 3.5 Flash generates a claim-level structured answer whose claims may only reference provision identifiers from the retrieved context. A six-layer verifier (schema, citation ID, temporal, numeric grounding, claim support, evidence completeness) decides between a verified answer and abstention with a standard reason code. The system abstains when evidence is insufficient instead of relying on unrestricted web search.

The study proposes a reviewed gold set of 200 questions split into 40 development, 40 validation, and 120 final test cases, together with four experiment suites: parser benchmark (P1-P3), embedding benchmark (E1-E3), retrieval ablation (R1-R10), and generation/verification ablation (G1-G7). Primary metrics include Recall@k, MRR@10, nDCG@10, Citation Precision/Recall/F1, Temporal Validity Accuracy, Numeric Grounding Accuracy, Evidence Completeness, and Abstention F1. Langfuse provides observability outside the correctness-critical path; RAGFlow is used only as an external comparison baseline.

\draftnotice

\noindent\textbf{Keywords:} Retrieval-Augmented Generation, Vietnamese traffic law, temporal retrieval, hybrid retrieval, citation verification, evidence completeness, Parser Router, Canonical Document IR.
\clearpage

\tableofcontents
\clearpage
\listoftables
\addcontentsline{toc}{chapter}{Danh mục bảng}
\clearpage
\listoffigures
\addcontentsline{toc}{chapter}{Danh mục hình}
\clearpage

\chapter*{Danh Mục Ký Hiệu và Từ Viết Tắt}
\addcontentsline{toc}{chapter}{Danh mục ký hiệu và từ viết tắt}

\begin{longtable}{p{3.2cm}p{10.2cm}}
\textbf{Ký hiệu} & \textbf{Diễn giải}\\
\toprule
AI & Artificial Intelligence - Trí tuệ nhân tạo\\
API & Application Programming Interface\\
BM25 & Best Matching 25 - thuật toán xếp hạng truy xuất từ khóa\\
Canonical IR & Canonical Document IR - biểu diễn tài liệu trung gian parser-neutral của dự án\\
CI/CD & Continuous Integration / Continuous Delivery\\
CRUD & Create, Read, Update, Delete\\
Dramatiq & Thư viện background processing cho Python, chạy worker ingestion\\
E2E & End-to-End\\
Evidence Gate & Evidence Completeness Gate - cổng kiểm tra tính đầy đủ bằng chứng trước khi sinh câu trả lời\\
F1 & F1-score - trung bình điều hòa của Precision và Recall\\
HITL & Human-in-the-Loop - con người trong vòng lặp (ở khâu review ingestion)\\
ID & Identifier - định danh\\
IR & Intermediate Representation - biểu diễn trung gian (Canonical Document IR)\\
JSON & JavaScript Object Notation\\
Langfuse & Nền tảng observability, quản lý prompt và experiment cho LLM\\
LegalProvision & Đơn vị pháp lý (Điều, Khoản, Điểm) kèm khoảng hiệu lực và provenance\\
LLM & Large Language Model - mô hình ngôn ngữ lớn\\
MinIO & Object storage S3-compatible\\
MRR & Mean Reciprocal Rank\\
nDCG & Normalized Discounted Cumulative Gain\\
OCR & Optical Character Recognition\\
P0/P1 & Mức ưu tiên bắt buộc / mở rộng\\
Parser Router & Bộ định tuyến parser: Docling chính, MinerU phụ/fallback\\
PII & Personally Identifiable Information\\
RAG & Retrieval-Augmented Generation\\
RAGFlow & Công cụ RAG mã nguồn mở, chỉ dùng làm baseline so sánh bên ngoài\\
RRF & Reciprocal Rank Fusion\\
SDK & Software Development Kit\\
SHA & Secure Hash Algorithm\\
SQL & Structured Query Language\\
SSE & Server-Sent Events\\
VLM & Vision-Language Model\\
\bottomrule
\end{longtable}
\clearpage

\chapter*{Nhận Xét Của Giảng Viên Hướng Dẫn}
\addcontentsline{toc}{chapter}{Nhận xét của giảng viên hướng dẫn}

\vspace{1cm}
\noindent\rule{\textwidth}{0.4pt}\\[1.2cm]
\noindent\rule{\textwidth}{0.4pt}\\[1.2cm]
\noindent\rule{\textwidth}{0.4pt}\\[1.2cm]
\noindent\rule{\textwidth}{0.4pt}\\[1.2cm]
\noindent\rule{\textwidth}{0.4pt}\\[1.2cm]
\noindent\rule{\textwidth}{0.4pt}\\[1.2cm]
\noindent\rule{\textwidth}{0.4pt}
\clearpage

% ============================================================
% MAIN CONTENT
% ============================================================
\pagenumbering{arabic}
\setcounter{page}{1}

% ============================================================
\chapter{Tổng Quan Đề Tài}
% ============================================================

\section{Bối cảnh}

Văn bản pháp luật giao thông đường bộ có ảnh hưởng trực tiếp đến hoạt động hằng ngày của người dân. Tuy nhiên, việc tra cứu thường đòi hỏi người dùng xác định đúng loại văn bản, số hiệu, Điều, Khoản, Điểm và trạng thái hiệu lực. Một quy định về mức phạt có thể được ban hành trong một Nghị định, được sửa đổi bởi văn bản khác và được thay thế ở một thời điểm cụ thể. Vì vậy, câu trả lời đúng ở năm 2023 chưa chắc đúng ở năm 2026.

Cơ sở dữ liệu quốc gia về văn bản pháp luật cung cấp toàn văn, thuộc tính, lịch sử và văn bản liên quan. Ví dụ, Luật số 36/2024/QH15 có hiệu lực từ ngày 01/01/2025 \cite{luat36_2024}; Nghị định số 168/2024/NĐ-CP cũng có hiệu lực từ ngày 01/01/2025 và quy định xử phạt vi phạm hành chính, trừ điểm và phục hồi điểm giấy phép lái xe trong lĩnh vực giao thông đường bộ \cite{nghidinh168_2024}. Trước thời điểm đó, nhiều tình huống xử phạt được tra cứu theo Nghị định số 100/2019/NĐ-CP và các văn bản sửa đổi \cite{nghidinh100_2019}. Điều này làm cho thời gian hiệu lực và quan hệ giữa các văn bản trở thành dữ liệu bắt buộc thay vì metadata phụ.

Các mô hình ngôn ngữ lớn có khả năng diễn đạt và tổng hợp tốt, nhưng không bảo đảm tự động rằng số hiệu hoặc trích dẫn được sinh ra là có thật. LegalCiteBench cho thấy hướng dẫn bằng prompt về sự không chắc chắn có thể giảm một phần câu trả lời tự tin sai, nhưng không tự cải thiện độ đúng của citation khi thiếu grounding bên ngoài \cite{chen2026legalcitebench}. Do đó, một hệ thống pháp luật cần ràng buộc mô hình bằng dữ liệu truy xuất và một lớp kiểm chứng xác định.

\section{Vấn đề nghiên cứu}

Một pipeline RAG cơ bản thường gồm bốn bước: chuyển tài liệu thành text, chia thành chunk, truy xuất chunk gần với câu hỏi và yêu cầu LLM tạo câu trả lời. Pipeline này còn bảy hạn chế khi áp dụng cho pháp luật:

\begin{enumerate}
  \item \textbf{Mất cấu trúc pháp lý:} fixed-token chunking có thể tách phần dẫn của Khoản khỏi các Điểm bên dưới hoặc gộp nhiều điều khoản không liên quan. Ranh giới pháp lý phải trùng ranh giới trích dẫn; nhãn Điểm tiếng Việt (kể cả \code{đ)}) không thể được xử lý bằng giả định bảng chữ cái đơn giản.
  \item \textbf{Không nhận biết thời gian:} vector search chỉ đo độ gần ngữ nghĩa, không biết văn bản đã hết hiệu lực hay chưa có hiệu lực tại ngày được hỏi.
  \item \textbf{Bỏ sót theo định danh chính xác:} retrieval thuần vector có thể bỏ sót số hiệu văn bản, số Điều, số Khoản hoặc cụm từ pháp lý có tính quyết định.
  \item \textbf{Quan hệ tham chiếu chéo chưa được mô hình hóa:} một Điều dẫn chiếu Điều khác, quy định xử phạt thường gắn kèm quy định trừ điểm; hệ thống chỉ trả lời được nửa thông tin nếu không mở rộng ngữ cảnh theo quan hệ.
  \item \textbf{Câu hỏi đa bằng chứng:} một câu hỏi có thể yêu cầu vừa mức phạt vừa số điểm trừ; hệ thống không được âm thầm trả lời một nửa dễ của câu hỏi.
  \item \textbf{Citation không kiểm chứng:} mô hình có thể gõ số Điều hoặc số hiệu không tồn tại, hoặc ghép nội dung đúng với \code{provision\_id} sai. Một citation tồn tại chưa đủ; số liệu trong claim còn phải khớp bằng chứng.
  \item \textbf{Đánh giá không tách lớp:} một điểm tổng hợp không chỉ ra lỗi đến từ parser, extraction, retrieval, temporal filter, generation hay verification.
\end{enumerate}

Vấn đề trung tâm của đề tài được phát biểu như sau:

\begin{quote}
Làm thế nào để xây dựng một hệ thống RAG truy xuất đúng đơn vị pháp lý, áp dụng đúng phiên bản văn bản tại thời điểm được hỏi, mở rộng đúng các quy định liên quan và chỉ trả về các nhận định có trích dẫn kiểm chứng được?
\end{quote}

\section{Mục tiêu}

\subsection{Mục tiêu tổng quát}

Xây dựng và đánh giá một hệ thống RAG nhận biết cấu trúc, quan hệ tham chiếu và thời gian hiệu lực để hỗ trợ tra cứu pháp luật giao thông đường bộ Việt Nam bằng tiếng Việt.

\subsection{Mục tiêu cụ thể}

\begin{enumerate}
  \item Xây dựng Parser Router định tuyến tài liệu qua Docling hoặc MinerU theo đặc tính tài liệu và quality gate.
  \item Thiết kế Canonical Document IR parser-neutral làm biểu diễn trung gian cho toàn bộ phân tích pháp lý.
  \item Xây dựng Legal Structure Extractor nhận diện phân cấp Chương, Mục, Điều, Khoản, Điểm với nhãn tiếng Việt a) b) c) d) đ) e) và giữ lại Điểm ngắn hợp lệ.
  \item Mô hình hóa quan hệ tham chiếu chéo và quan hệ hiệu lực bằng PostgreSQL, không dùng Neo4j.
  \item \begin{sloppypar} Biểu diễn văn bản và điều khoản bằng LegalDocument và LegalProvision, kèm định danh ổn định và khoảng hiệu lực \code{[effective\_from, effective\_to)}. \end{sloppypar}
  \item Kết hợp exact lookup, dense retrieval và sparse BM25 trong Qdrant bằng RRF, kèm reranking và mở rộng ngữ cảnh pháp lý.
  \item Xây dựng evidence plan và Evidence Completeness Gate để không trả lời nửa vời cho câu hỏi đa bằng chứng.
  \item Sinh câu trả lời có cấu trúc cấp claim và kiểm chứng citation theo \code{provision\_id} ở sáu tầng.
  \item Thiết kế abstention có lý do chuẩn khi thiếu dữ liệu hoặc verification thất bại.
  \item Xây dựng gold set 200 câu và chạy bốn suite thí nghiệm A-D.
  \item Thiết lập observability qua Langfuse ngoài đường tới hạn và baseline so sánh RAGFlow bên ngoài.
\end{enumerate}

\section{Câu hỏi nghiên cứu}

\begin{table}[H]
\caption{Câu hỏi nghiên cứu}
\begin{tabularx}{\textwidth}{p{1.2cm}X}
\toprule
\textbf{Mã} & \textbf{Câu hỏi}\\
\midrule
RQ1 & Docling, MinerU và Parser Router khác nhau thế nào về chất lượng trích xuất cấu trúc pháp luật Việt Nam (Article/Clause/Point P/R/F1, Short Point Recall, Vietnamese đ) Recall, Parent Context Completeness, Provenance Coverage)?\\
RQ2 & Embedding nào trong ba ứng viên (Gemini Embedding 2, Jina Embeddings v5 text-nano, Jina Embeddings v5 text-small) cho kết quả retrieval tốt nhất trên câu hỏi pháp luật tiếng Việt (Recall@10, MRR@10, nDCG@10)?\\
RQ3 & Từng tầng của retrieval đa tầng (exact lookup, sparse/RRF, query normalization, multi-query rewrite, conditional HyDE, reranker, parent/sibling expansion, cross-reference expansion, temporal filter) cải thiện Recall@k, MRR và nDCG ở mức nào?\\
RQ4 & Temporal filtering làm giảm tỷ lệ trích dẫn sai phiên bản (Temporal Leakage Rate) ở mức nào và giữ được bao nhiêu trích dẫn hợp lệ cho câu hỏi lịch sử?\\
RQ5 & Verification sáu tầng, đặc biệt numeric grounding, claim support và evidence completeness, làm giảm invalid citation và unsupported claim ở mức nào, đồng thời ảnh hưởng thế nào đến tỷ lệ abstention?\\
RQ6 & Những nhóm lỗi nào giới hạn chất lượng hệ thống: parser, extraction, retrieval, temporal metadata, generation hay verification?\\
\bottomrule
\end{tabularx}
\end{table}

\section{Đối tượng và phạm vi}

Đối tượng nghiên cứu là pipeline RAG phục vụ hỏi đáp pháp luật giao thông đường bộ Việt Nam. Corpus mục tiêu gồm khoảng 20 đến 30 văn bản chính thống, có cả văn bản hiện hành, lịch sử và ít nhất 5 chuỗi sửa đổi, thay thế hoặc bãi bỏ. Người dùng cuối có thể hỏi quy định hiện hành, hỏi tại ngày cụ thể hoặc so sánh hai thời điểm.

Đề tài không bao phủ toàn bộ hệ thống pháp luật, không cung cấp kết luận thay thế cơ quan có thẩm quyền, không tìm web tự do để sinh answer, không xây dựng multi-agent, không fine-tune LLM và không triển khai ứng dụng di động hoặc voice.

\begin{table}[H]
\caption{Phân định phạm vi}
\begin{tabularx}{\textwidth}{X X}
\toprule
\textbf{Trong phạm vi} & \textbf{Ngoài phạm vi}\\
\midrule
PDF pháp luật giao thông đường bộ & Toàn bộ pháp luật Việt Nam\\
Parser Router: Docling chính, MinerU phụ/fallback & Crawler tự động activate tài liệu\\
Canonical Document IR và Legal Structure Extractor & Tư vấn pháp lý cá nhân hóa có tính kết luận\\
Quan hệ tham chiếu chéo trong PostgreSQL & Open web search làm nguồn answer\\
Temporal retrieval và sửa đổi từng phần & Autonomous multi-agent\\
Exact lookup + dense + sparse + RRF + reranker & Fine-tuning LLM\\
Evidence planning và Evidence Completeness Gate & Knowledge graph hoặc Neo4j\\
Verification sáu tầng và abstention & Mobile, voice, Kubernetes\\
Langfuse observability ngoài đường tới hạn & RAGFlow làm nền tảng chính\\
Gold set 200 câu và bốn suite A-D & Local LLM\\
RAGFlow làm baseline so sánh bên ngoài & Dịch vụ microservices\\
\bottomrule
\end{tabularx}
\end{table}

\section{Phương pháp thực hiện}

Nghiên cứu sử dụng phương pháp thiết kế và thực nghiệm hệ thống. Trước hết, yêu cầu và invariant pháp lý được phân tích để xây dựng data model. Tiếp theo, pipeline ingestion (Parser Router, Canonical Document IR, Legal Structure Extractor, resolver), retrieval đa tầng và verification được cài đặt theo kiến trúc module. Sau đó, các biến thể được đánh giá trên cùng corpus và gold set qua bốn suite thí nghiệm A-D. Kết quả được phân tích theo metric xác định, metric LLM phụ, latency, cost và error taxonomy.

Quá trình thực hiện trải qua một lần thay đổi thiết kế: phương án thiết kế đầu tiên dựa trên một framework trích xuất dữ liệu phi cấu trúc tách rời \cite{udef2026}. Sau khi đánh giá, phương án này được thay bằng Parser Router, Canonical Document IR và Legal Structure Extractor do dự án sở hữu để giảm tầng chuyển đổi trung gian và kiểm soát trực tiếp việc nhận diện phân cấp pháp luật Việt Nam. Lý do chi tiết được trình bày trong Chương 2.

\section{Đóng góp dự kiến}

Sáu đóng góp chính gồm:

\begin{enumerate}
  \item Parser Router định tuyến Docling/MinerU theo đặc tính tài liệu và quality gate, cùng Canonical Document IR parser-neutral do dự án sở hữu;
  \item Legal Structure Extractor cho phân cấp pháp luật Việt Nam, gồm nhãn Điểm tiếng Việt \code{d)} và \code{đ)}, short-Point retention và stable ID tái tạo được;
  \item mô hình quan hệ tham chiếu chéo cấp provision và cấp văn bản lưu trong PostgreSQL kèm mở rộng ngữ cảnh pháp lý theo quan hệ;
  \item retrieval đa tầng với evidence planning, Evidence Completeness Gate và mở rộng ngữ cảnh có ghi lý do (\code{added\_by}, \code{depth});
  \item verification xác định sáu tầng với bất biến Returned Invalid Citation Rate bằng 0;
  \item gold set tiếng Việt 200 câu và phương pháp evaluation bốn suite A-D có thể tái lập, kèm baseline RAGFlow bên ngoài.
\end{enumerate}

\section{Bố cục khóa luận}

Chương 2 trình bày nền tảng lý thuyết và nghiên cứu liên quan. Chương 3 mô tả yêu cầu, corpus, gold set và phương pháp thực nghiệm. Chương 4 trình bày thiết kế hệ thống. Chương 5 mô tả cài đặt và triển khai. Chương 6 trình bày kiểm thử và đánh giá. Chương 7 tổng kết, nêu hạn chế và hướng phát triển.

% ============================================================
\chapter{Cơ Sở Lý Thuyết và Nghiên Cứu Liên Quan}
% ============================================================

\section{Mô hình ngôn ngữ lớn}

Transformer sử dụng cơ chế attention để mô hình hóa quan hệ giữa các token trong chuỗi và là nền tảng của nhiều mô hình ngôn ngữ hiện đại \cite{vaswani2017attention}. LLM có khả năng sinh văn bản, phân loại và trích xuất có cấu trúc, nhưng đầu ra không mặc nhiên là bằng chứng. Trong hệ thống này, LLM chỉ đảm nhiệm phân tích phần ngôn ngữ chưa xử lý được bằng rule và sinh answer từ context đã giới hạn.

Generator được cấu hình theo structured output. Gemini 3.5 Flash hỗ trợ structured outputs theo \code{json\_schema} tương thích Pydantic và context lớn \cite{google2026gemini35flash}. Tuy nhiên, model capability không thay thế retrieval correctness hoặc citation verification.

\section{Retrieval-Augmented Generation}

RAG kết hợp retriever với generator để cung cấp external knowledge cho mô hình \cite{lewis2020rag}. Một hệ thống RAG có thể chia thành:

\begin{enumerate}
  \item ingestion và indexing;
  \item query processing;
  \item retrieval;
  \item context construction;
  \item generation;
  \item verification và response rendering.
\end{enumerate}

Đề tài xem RAG như một hệ thống nhiều lớp. Parser sai làm extraction sai; extraction sai làm index sai; retrieval sai đặt trần cho generation; generation đúng nhưng citation sai vẫn không thể trả ra người dùng. Vì vậy, một pipeline RAG tổng quát (PDF, chunk, vector, LLM) là chưa đủ cho pháp luật: nó không biết ranh giới trích dẫn, không biết phiên bản hiệu lực, không biết quan hệ giữa các quy định và không kiểm chứng được câu trả lời.

\section{Dense retrieval và embedding}

Dense retrieval ánh xạ query và document sang vector thực. Độ tương đồng cosine thường được dùng để xếp hạng. Gemini Embedding 2 hỗ trợ output dimension linh hoạt từ 128 đến 3072 và khuyến nghị các mức 768, 1536 hoặc 3072 \cite{google2026embedding2}. Hệ thống dùng 768 chiều làm cấu hình thử nghiệm khởi điểm để cân bằng chất lượng và chi phí lưu trữ, nhưng embedding production chưa được chốt vĩnh viễn: Suite B benchmark ba ứng viên (Gemini Embedding 2, Jina Embeddings v5 text-nano, Jina Embeddings v5 text-small) trên câu hỏi pháp luật tiếng Việt trước khi chọn.

Embedding text của một provision không chỉ chứa nội dung gốc mà còn có tên văn bản, số hiệu và parent context (\code{retrieval\_text}). Việc bổ sung context giúp các Điểm ngắn vẫn mang đủ nghĩa khi được truy xuất độc lập, trong khi \code{source\_text} luôn giữ nguyên văn bản pháp lý gốc.

\section{Sparse retrieval và BM25}

BM25 là mô hình xác suất dựa trên tần suất từ, độ dài tài liệu và độ hiếm của thuật ngữ \cite{robertson2009bm25}. Một dạng công thức phổ biến là:

\begin{equation}
\operatorname{BM25}(q,d)=\sum_{t\in q}\operatorname{IDF}(t)
\frac{f(t,d)(k_1+1)}{f(t,d)+k_1\left(1-b+b\frac{|d|}{\operatorname{avgdl}}\right)}.
\end{equation}

Sparse retrieval đặc biệt hữu ích cho số hiệu, tên văn bản, cụm từ pháp lý và tham chiếu Điều/Khoản/Điểm. Hệ thống không loại mạnh các từ như ``không'', ``phải'', ``được'', ``trừ'', vì chúng có thể thay đổi nghĩa quy phạm. Qdrant cung cấp sparse/BM25 retrieval ngay trong collection bằng tokenizer-based model \cite{qdrant2026bm25}; với tiếng Việt cần kiểm chứng tokenizer trong Suite C trước khi khóa cấu hình.

\section{Hybrid retrieval và RRF}

Dense và sparse score nằm trên scale khác nhau. Reciprocal Rank Fusion kết hợp các danh sách theo thứ hạng thay vì raw score \cite{cormack2009rrf}. Với document $d$ và tập ranking $R$:

\begin{equation}
\operatorname{RRF}(d)=\sum_{r\in R}\frac{1}{k+\operatorname{rank}_r(d)}.
\end{equation}

Qdrant hỗ trợ named dense/sparse vectors, prefetch và RRF trong Query API \cite{qdrant2026hybrid}. Điều này loại bỏ nhu cầu duy trì hai hệ retrieval tách rời và code fusion thủ công.

\section{Chunking tài liệu}

Chunking quyết định đơn vị được embed và retrieve. Fixed-token chunking đơn giản nhưng có thể phá cấu trúc. Docling cung cấp HierarchicalChunker và HybridChunker; HybridChunker thực hiện split chunk quá dài và merge peer chunk nhỏ dựa trên tokenizer, đồng thời giữ metadata của document structure \cite{docling2026chunking}.

Đối với pháp luật, hierarchy tổng quát chưa đủ. Một Điểm có thể phụ thuộc vào lead-in của Khoản; một Khoản có thể phụ thuộc tiêu đề Điều. Ranh giới chunk phải trùng ranh giới trích dẫn. Vì vậy, Legal Structure Extractor của dự án xác định ranh giới pháp lý trước, còn chunking chỉ xử lý provision quá dài theo câu nhưng giữ cùng \code{provision\_id}; Docling chunker không được dùng làm legal structure parser chính.

\begin{figure}[H]
\centering
\begin{tikzpicture}[
  node distance=8mm,
  box/.style={draw, rounded corners, minimum width=5.2cm, minimum height=0.75cm, align=center, fill=SoftBlue},
  arrow/.style={-{Latex[length=2mm]}, thick}
]
\node[box] (doc) {Văn bản pháp luật};
\node[box, below=of doc] (chapter) {Chương};
\node[box, below=of chapter] (section) {Mục};
\node[box, below=of section] (article) {Điều};
\node[box, below=of article] (clause) {Khoản};
\node[box, below=of clause] (point) {Điểm (a) b) c) d) đ) e))};
\draw[arrow] (doc) -- (chapter);
\draw[arrow] (chapter) -- (section);
\draw[arrow] (section) -- (article);
\draw[arrow] (article) -- (clause);
\draw[arrow] (clause) -- (point);
\end{tikzpicture}
\caption{Cấu trúc phân cấp của một đơn vị pháp lý}
\label{fig:legal-hierarchy}
\end{figure}

\section{Document parsing và provenance}

Docling chuyển nhiều định dạng tài liệu sang một representation giàu cấu trúc, nhận biết layout và table, đồng thời hỗ trợ OCR tiếng Việt qua Tesseract, EasyOCR hoặc RapidOCR \cite{auer2024docling}. MinerU là parser phụ và fallback/challenger với OCR 109 ngôn ngữ và backend pipeline chạy CPU. Các khả năng chi tiết của hai parser được mô tả trong tài liệu nghiên cứu công nghệ của dự án; không khẳng định parser nào vượt trội tuyệt đối cho mọi trường hợp.

Hệ thống không đưa đầu ra parser trực tiếp vào phân tích pháp lý. Parser Router chọn parser theo đặc tính tài liệu và quality gate, sau đó đầu ra được chuyển sang Canonical Document IR do dự án sở hữu gồm \code{ParsedDocument}, \code{ParsedPage} và \code{DocumentElement}. Tầng IR cô lập Legal Structure Extractor khỏi định dạng đầu ra của Docling/MinerU; thay hoặc nâng cấp parser chỉ cần một adapter mới.

\begin{sloppypar}
Provenance được giữ dưới dạng page number, bounding box và \code{source\_element\_ids} truy vết về IR. Provenance phục vụ ba mục tiêu: review extraction, hiển thị đoạn nguồn và audit khi answer bị phản ánh. Một citation không chỉ trỏ đến văn bản mà còn phải mở được passage cụ thể.
\end{sloppypar}

\section{Temporal data và phiên bản pháp luật}

Một provision hợp lệ tại ngày $d$ khi:

\begin{equation}
\begin{aligned}
\operatorname{valid}(p,d)=
&\left(p.\operatorname{effective\_from}\le d\right)\\
&\land\ \left(p.\operatorname{effective\_to}=\varnothing\ \lor\ d<p.\operatorname{effective\_to}\right)\\
&\land\ \left(p.\operatorname{review\_status}=\text{ACCEPTED}\right).
\end{aligned}
\end{equation}

Khoảng thời gian được biểu diễn theo dạng nửa kín $[\operatorname{from},\operatorname{to})$. Khi một version mới có hiệu lực, version cũ kết thúc tại chính ngày bắt đầu của version mới. Partial amendment được xử lý ở mức provision thay vì đóng toàn bộ document: \code{provision\_id} giữ nguyên, nội dung mới được lưu dưới version mới với khoảng hiệu lực mới. Các sự kiện pháp lý (có hiệu lực, sửa đổi, thay thế, bãi bỏ, đính chính) được ghi nhận bằng \code{LegalEffectEvent} để phục vụ temporal resolver và câu hỏi so sánh lịch sử.

Thời gian hiệu lực không phải metadata phụ mà là điều kiện lọc bắt buộc của mọi retrieval: câu trả lời đúng cho một hành vi năm 2023 phải dùng văn bản hợp lệ tại năm 2023, kể cả khi văn bản đó đã bị thay thế. Hệ thống không dùng văn bản hiện hành làm mặc định cho câu hỏi có thể là lịch sử.

\section{Citation verification và abstention}

LegalBench-RAG nhấn mạnh retrieval chính xác ở các đoạn pháp lý nhỏ thay vì trả cả document dài \cite{pipitone2024legalbenchrag}. Legal RAG Bench tiếp tục chỉ ra retrieval là yếu tố chủ đạo trong hiệu năng end-to-end của legal RAG \cite{butler2026legalragbench}. LegalCiteBench cho thấy citation generation không grounding là một failure mode nghiêm trọng \cite{chen2026legalcitebench}.

Từ các kết quả này, hệ thống không yêu cầu LLM viết citation text. Mỗi claim chỉ được chọn \code{provision\_id} từ whitelist context. Tuy nhiên, citation tồn tại là chưa đủ: một trích dẫn hợp lệ về định danh vẫn có thể đi kèm số liệu sai (mức phạt, số điểm trừ) hoặc claim không được passage hỗ trợ. Vì vậy, ngoài kiểm tra ID và hiệu lực, hệ thống còn kiểm tra numeric grounding và claim support. Abstention được coi là hành vi đúng khi evidence không đủ.

\section{Evidence completeness}

Nhiều câu hỏi pháp luật yêu cầu nhiều loại bằng chứng phối hợp. Ví dụ, câu hỏi về mức phạt và điểm trừ giấy phép lái xe cần cả \code{monetary\_penalty} và \code{license\_points}. Nếu pipeline chỉ truy xuất được mức phạt và trả lời luôn, người dùng nhận được câu trả lời nửa vời. Đề tài xử lý vấn đề này bằng evidence plan (danh sách loại bằng chứng cần thiết được tạo trong query understanding) và Evidence Completeness Gate trước khi sinh câu trả lời. Khi còn thiếu loại bằng chứng bắt buộc, hệ thống chạy targeted retrieval hoặc mở rộng ngữ cảnh theo quan hệ thay vì sinh câu trả lời.

Bằng chứng cho câu trả lời thường nằm rải rác ở nhiều provision liên quan. Mô hình quan hệ tham chiếu chéo (trình bày trong Chương 4) là cơ sở để mở rộng ngữ cảnh quanh các seed provision mạnh, đồng thời là dữ liệu gold cho metric Cross-reference Resolution Recall và Multi-hop Evidence Completeness.

\section{Đánh giá hệ thống RAG}

RAGAS cung cấp các metric như faithfulness, context precision và context recall \cite{es2024ragas,ragas2026metrics}. Tuy nhiên, metric LLM có thể biến động theo judge và prompt. Vì vậy, khóa luận dùng metric xác định làm headline (Recall@k, MRR, nDCG, Citation P/R/F1, Temporal Validity Accuracy, Numeric Grounding Accuracy, Evidence Set Recall, Abstention P/R/F1), còn Ragas và judge độc lập GPT-5.4 mini chỉ bổ sung đánh giá semantic \cite{openai2026gpt54mini}.

\section{Khoảng trống nghiên cứu}

Các benchmark pháp luật hiện có chủ yếu dùng tiếng Anh, hợp đồng hoặc án lệ. Đề tài tập trung vào văn bản quy phạm pháp luật giao thông Việt Nam, nơi hierarchy, nhãn Điểm tiếng Việt và hiệu lực theo thời gian là những đặc điểm bắt buộc. Khoảng trống cụ thể gồm:

\begin{itemize}
  \item chưa có benchmark nhỏ, review được cho current/historical/comparison query tiếng Việt;
  \item pháp luật Việt Nam chưa được đánh giá trên sự khác biệt giữa parser tổng quát và parser pháp lý chuyên biệt;
  \item citation ID, temporal validity và numeric grounding thường chưa được kiểm chứng cùng lúc;
  \item nhiều pipeline trộn web fallback vào answer path, làm giảm khả năng tái lập;
  \item câu hỏi đa bằng chứng và quan hệ tham chiếu chéo chưa được đánh giá như hạng mục đầu tiên trong legal RAG benchmark.
\end{itemize}

\begin{table}[H]
\caption{So sánh định hướng nghiên cứu liên quan}
\small
\begin{tabularx}{\textwidth}{p{2.7cm}p{1.8cm}p{2.0cm}p{1.6cm}Y}
\toprule
\textbf{Công trình} & \textbf{Retrieval} & \textbf{Citation} & \textbf{Temporal} & \textbf{Định hướng chính}\\
\midrule
RAGAS & Có metric & Gián tiếp & Không & Đánh giá RAG tổng quát\\
LegalBench-RAG & Chính xác & Không là trọng tâm & Không & Retrieval đoạn pháp lý nhỏ\\
Legal RAG Bench & End-to-end & Độ bám nguồn & Hạn chế & Phân rã lỗi retrieval/reasoning\\
LegalCiteBench & Closed-book & Trọng tâm & Không & Citation recovery và verification\\
Đề tài này & Exact + dense + sparse & ID + sáu tầng verifier & Có & Luật giao thông Việt Nam\\
\bottomrule
\end{tabularx}
\end{table}

% ============================================================
\chapter{Phân Tích Yêu Cầu và Phương Pháp Nghiên Cứu}
% ============================================================

\section{Tác nhân và use case}

Ba tác nhân chính là người dùng cuối, corpus reviewer và nhà phát triển. Người dùng hỏi luật, xem citation và gửi feedback. Reviewer kiểm tra metadata, hierarchy, provenance, quan hệ và hiệu lực trước khi index. Nhà phát triển quản lý corpus, model, deployment, evaluation và baseline RAGFlow. HITL chỉ nằm ở khâu review trong ingestion pipeline, không nằm trong online query.

\begin{figure}[H]
\centering
\resizebox{0.94\textwidth}{!}{%
\begin{tikzpicture}[
  actor/.style={draw, rounded corners, minimum width=3.2cm, minimum height=0.9cm, fill=SoftGray, align=center},
  system/.style={draw, rounded corners, minimum width=6.0cm, minimum height=1.0cm, fill=SoftBlue, align=center},
  arrow/.style={-{Latex[length=2mm]}, thick}
]
\node[system] (sys) {Structure-Aware and Temporal RAG};
\node[actor, left=2.2cm of sys] (user) {Người dùng cuối};
\node[actor, above=1.2cm of sys] (reviewer) {Corpus Reviewer};
\node[actor, right=2.2cm of sys] (dev) {Nhà phát triển};
\draw[arrow] (user) -- node[above, font=\small]{Câu hỏi / citation / feedback} (sys);
\draw[arrow] (reviewer) -- node[right, font=\small]{Review ingestion} (sys);
\draw[arrow] (dev) -- node[above, font=\small]{Config / evaluation} (sys);
\end{tikzpicture}%
}
\caption{Các tác nhân chính của hệ thống}
\end{figure}

Hệ thống có mười hai use case: UC-01 hỏi luật hiện hành, UC-02 hỏi luật lịch sử, UC-03 so sánh quy định, UC-04 tìm provision, UC-05 xem citation và nguồn, UC-06 abstention, UC-07 ingest tài liệu bất đồng bộ, UC-08 review ingestion, UC-09 chạy evaluation, UC-10 gửi feedback, UC-11 chạy parser benchmark (Suite A) và UC-12 corpus QA review.

\section{Yêu cầu chức năng}

\begin{longtable}{p{1.6cm}p{5.2cm}p{6.5cm}}
\caption{Yêu cầu chức năng P0 (rút gọn)}\\
\toprule
\textbf{Mã} & \textbf{Yêu cầu} & \textbf{Tiêu chí kiểm chứng}\\
\midrule
\endfirsthead
\toprule
\textbf{Mã} & \textbf{Yêu cầu} & \textbf{Tiêu chí kiểm chứng}\\
\midrule
\endhead
FR-01 & Parser Router (Docling | MinerU) & Routing đúng theo đặc tính tài liệu; quality gate kích hoạt fallback\\
FR-02 & Canonical Document IR parser-neutral & Extractor chỉ đọc IR; adapter mới không đổi extractor\\
FR-03 & Legal Structure Extractor & Phân cấp đúng; nhãn \code{d)}/\code{đ)} không lẫn; giữ Điểm ngắn\\
FR-04 & Parent-context enrichment & \code{source\_text} bất biến; \code{retrieval\_text} kế thừa context cha\\
FR-05 & Legal Reference Resolver & Quan hệ khớp fixture; unresolved định tuyến review\\
FR-06 & Temporal and Amendment Resolver & Interval đúng quy tắc; hiệu lực không chắc chắn định tuyến review\\
FR-07 & Background ingestion qua hàng đợi & Upload trả 202; actor idempotent; không parse đồng bộ\\
FR-08 & Object storage qua MinIO & Round-trip put/get; backup sang nơi độc lập\\
FR-09 & Review routing trước khi index & Chỉ accepted được index; dropped không bao giờ index\\
FR-10 & Corpus QA & Báo cáo đủ 16 chỉ số\\
FR-11 & Query Understanding và evidence plan & Intent/date/refs đúng; evidence plan liệt kê đủ loại bằng chứng\\
FR-12 & Query Expansion & Giữ câu gốc; rewrite có giới hạn; HyDE chỉ bật có điều kiện\\
FR-13 & Exact legal lookup & Trả đúng provision theo số hiệu, Điều, Khoản, Điểm\\
FR-14 & Dense + sparse + RRF & Ba kênh chạy độc lập; variant tái lập bằng config\\
FR-15 & Reranking & Chạy như stage chuẩn; không khẳng định cải thiện trước benchmark\\
FR-16 & Legal context expansion & Chỉ mở rộng quanh seed mạnh; ghi \code{added\_by} và \code{depth}\\
FR-17 & Evidence Completeness Gate & Câu đa bằng chứng thiếu loại bằng chứng bị đánh dấu INCOMPLETE\\
FR-18 & Hỏi luật hiện hành & Chỉ cite provision có hiệu lực tại ngày hỏi\\
FR-19 & Hỏi luật lịch sử & Dùng phiên bản hợp lệ tại mốc hỏi, không dùng hiện hành làm mặc định\\
FR-20 & So sánh hai thời điểm & Hai temporal contexts độc lập; không gộp citation hai mốc\\
FR-21 & Tìm kiếm provision & Trả top-k kèm hierarchy, hiệu lực và provenance; không gọi generator\\
FR-22 & Structured generation & Output đúng schema cấp claim; chỉ dùng ID trong whitelist\\
FR-23 & Verification sáu tầng L1-L6 & Invalid citation không trả ra; claim số liệu sai bị chặn\\
FR-24 & Failure-aware repair và abstention & Mỗi lỗi có đường sửa riêng; hết giới hạn thì ABSTAIN\\
FR-25 & Disclaimer & Có trong mọi verified/abstained response\\
FR-26 & Observability qua Langfuse & Trace ghi lại được; Langfuse down không fail query\\
FR-27 & Feedback người dùng cuối & Round-trip; gắn đúng \code{trace\_id}; danh mục báo cáo đầy đủ\\
FR-28 & Evaluation và ablation & Bốn suite A-D chạy được; run lưu đủ metadata\\
FR-31 & Benchmark đối chiếu RAGFlow & Bốn variant trên cùng corpus và eval queries; chạy riêng\\
FR-32 & Hiển thị answer và citation & Citation dựng từ metadata; không stream draft chưa verify\\
\bottomrule
\end{longtable}

Các yêu cầu FR-29 (lưu lịch sử truy vấn) và FR-30 (admin review UI) là P1, chỉ triển khai sau khi mọi acceptance criteria P0 đạt; trong P0, reviewer dùng CLI và JSON report.

\section{Yêu cầu phi chức năng}

Các invariant quan trọng gồm: returned invalid citation rate bằng 0; unverified answer không được trả; mọi indexed provision có page provenance và thuộc document đã accept; final evaluation phải lưu corpus hash và gold-set hash; hệ thống chạy local bằng Docker Compose với bảy service. Các ngưỡng hiệu năng là mục tiêu kỹ thuật trong môi trường test, không phải kết quả đo: retrieval P95 không quá 2 giây, end-to-end P50 không quá 12 giây, end-to-end P95 không quá 20 giây, search API P95 không quá 3 giây. Ingestion chạy nền với \code{MAX\_INGESTION\_WORKERS=1}. Langfuse không nằm trên đường tới hạn; nếu không khả dụng, query vẫn hoạt động.

\section{Thiết kế corpus}

Corpus được chọn theo bốn tiêu chí:

\begin{enumerate}
  \item nguồn chính thống và có thể kiểm tra lại;
  \item liên quan giao thông đường bộ;
  \item có metadata ban hành và hiệu lực;
  \item tạo được nhóm current, historical và relation chain.
\end{enumerate}

Quy mô mục tiêu là 20 đến 30 văn bản chính thống, có ít nhất 5 chuỗi sửa đổi, thay thế hoặc bãi bỏ, có cả văn bản hiện hành và lịch sử. Mỗi file có manifest chứa document ID, số hiệu, ngày ban hành, khoảng hiệu lực, URL nguồn, file hash, relation notes và review status. Ưu tiên PDF born-digital có text layer; tài liệu scan đi qua OCR backend CPU và phải qua quality gate trước khi review.

\begin{lstlisting}[language=json,caption={Ví dụ corpus manifest rút gọn}]
{
  "document_id": "nd-168-2024",
  "document_number": "168/2024/ND-CP",
  "document_type": "DECREE",
  "issued_date": "2024-12-26",
  "effective_from": "2025-01-01",
  "effective_to": null,
  "source_url": "https://vbpl.vn/...",
  "file_hash": "sha256:...",
  "relation_notes": "SUPERSEDES nd-100-2019",
  "review_status": "ACCEPTED"
}
\end{lstlisting}

Corpus có báo cáo/dashboard chất lượng riêng (corpus QA) với các chỉ số: document count, article count, clause count, point count, Point coverage, short-Point retention, tỷ lệ phát hiện nhãn \code{đ)}, orphan Point count, orphan Clause count, duplicate provision count, parent-context coverage, provenance coverage, table coverage, unresolved cross-reference count, unknown effective date count và temporal conflict count. Với các văn bản quan trọng (ví dụ Nghị định 168), thực hiện structural QA có mục tiêu riêng.

\section{Thiết kế gold set}

\begin{sloppypar}
Gold set mục tiêu gồm 200 câu đã review, chia 40 development (lặp phát triển), 40 validation (chọn ngưỡng/model/prompt) và 120 final test (báo cáo kết quả, đóng băng, không dùng để tuning). Có 17 danh mục bắt buộc: CURRENT, HISTORICAL, COMPARISON, EXACT\_REFERENCE, PENALTY, LICENSE\_POINTS, CONDITION, EXCEPTION, PROCEDURE, CROSS\_REFERENCE, MULTI\_PROVISION, MULTI\_DOCUMENT, COLLOQUIAL\_QUERY, AMBIGUOUS, MISSING\_INFORMATION, OUT\_OF\_SCOPE, ADVERSARIAL\_CITATION.
\end{sloppypar}

\begin{sloppypar}
Mỗi record chứa id, question, category, query\_date, comparison\_dates, vehicle\_type, expected\_provision\_ids, acceptable\_provision\_ids, required\_evidence, must\_include\_facts, must\_not\_include\_facts, temporal\_metadata, review\_status, reviewed\_by, gold\_version và hash. Mỗi câu có thể yêu cầu nhiều provision bằng chứng; \code{required\_evidence} liệt kê các loại bằng chứng như \code{violation\_definition}, \code{monetary\_penalty}, \code{license\_points}, \code{license\_suspension}, \code{exception}, \code{procedure}, \code{legal\_condition}.
\end{sloppypar}

\section{Thiết kế thí nghiệm}

Bốn suite thí nghiệm thay thế mô hình ablation đơn lẻ của thiết kế trước:

\begin{table}[H]
\caption{Suite A - Parser benchmark}
\small
\begin{tabularx}{\textwidth}{p{1.2cm}Y}
\toprule
\textbf{Thí nghiệm} & \textbf{Parser}\\
\midrule
P1 & Docling\\
P2 & MinerU\\
P3 & Parser Router (quyết định routing + quality gates)\\
\bottomrule
\end{tabularx}
\end{table}

Chỉ số Suite A: Article P/R/F1, Clause P/R/F1, Point P/R/F1, Short Point Recall, Vietnamese đ) Recall, Parent Context Completeness, Table Preservation, Header/Footer Leakage, Provenance Coverage.

\begin{table}[H]
\caption{Suite B - Embedding benchmark}
\small
\begin{tabularx}{\textwidth}{p{1.2cm}p{6.0cm}p{3.0cm}}
\toprule
\textbf{Thí nghiệm} & \textbf{Model} & \textbf{Dimensions}\\
\midrule
E1 & Gemini Embedding 2 & 768\\
E2 & Jina Embeddings v5 text-nano & 768\\
E3 & Jina Embeddings v5 text-small & 1024\\
\bottomrule
\end{tabularx}
\end{table}

Chỉ số Suite B: Recall@10, MRR@10, nDCG@10 trên câu hỏi pháp luật tiếng Việt, latency, cost.

\begin{table}[H]
\caption{Suite C - Retrieval ablation}
\small
\begin{tabularx}{\textwidth}{p{1.2cm}Y}
\toprule
\textbf{Thí nghiệm} & \textbf{Cấu hình}\\
\midrule
R1 & Legal chunk + dense\\
R2 & R1 + sparse/RRF\\
R3 & R2 + query normalization\\
R4 & R3 + multi-query rewrite\\
R5 & R4 + conditional HyDE\\
R6 & R5 + reranker\\
R7 & R6 + parent/sibling expansion\\
R8 & R7 + cross-reference expansion\\
R9 & R8 + temporal filtering\\
R10 & Complete retrieval pipeline\\
\bottomrule
\end{tabularx}
\end{table}

\begin{table}[H]
\caption{Suite D - Generation và verification ablation}
\small
\begin{tabularx}{\textwidth}{p{1.2cm}Y}
\toprule
\textbf{Thí nghiệm} & \textbf{Cấu hình}\\
\midrule
G1 & Prompt-only\\
G2 & Structured output\\
G3 & G2 + citation ID verifier\\
G4 & G3 + temporal verifier\\
G5 & G4 + numeric grounding\\
G6 & G5 + claim support\\
G7 & G6 + evidence completeness\\
\bottomrule
\end{tabularx}
\end{table}

Ma trận thí nghiệm được thiết kế có tham khảo các quan sát bên ngoài từ dự án Traffic-RAG (ranh giới chunk trùng ranh giới trích dẫn, Điểm ngắn không được lọc bỏ, \code{đ)} phải được nhận diện, retrieval cần ngữ cảnh câu mở đầu Khoản, mở rộng Khoản lân cận, citation filtering đơn thuần là chưa đủ, label gold có thể sai cần review độc lập). Các quan sát này chỉ là động lực thiết kế thí nghiệm, không phải kết quả của VNLRAG.

\section{Metric}

Với gold set $G_q$ và top-k result $R_q^k$:

\begin{equation}
\operatorname{Recall@k}(q)=\frac{|G_q\cap R_q^k|}{|G_q|}.
\end{equation}

Nếu relevant item đầu tiên ở rank $r_q$:

\begin{equation}
\operatorname{MRR}=\frac{1}{N}\sum_{q=1}^{N}\frac{1}{r_q}.
\end{equation}

Citation Precision, Recall và F1 được tính bằng exact/acceptable provision IDs. Temporal Validity Accuracy đo tỷ lệ citation hợp lệ tại query date; Temporal Leakage Rate đo tỷ lệ citation dùng phiên bản sai thời điểm. Abstention được đánh giá bằng precision, recall, F1 và confusion matrix.

\begin{table}[H]
\caption{Nhóm metric đánh giá}
\begin{tabularx}{\textwidth}{p{3.2cm}X}
\toprule
\textbf{Lớp} & \textbf{Metric}\\
\midrule
Corpus & Hierarchy F1, Point Coverage, Short Point Recall, Provenance Coverage, Parent Context Coverage\\
Retrieval & Recall@5/10/20, MRR@10, nDCG@10\\
Evidence & Evidence Set Recall, All Required Evidence@10, Cross-reference Resolution Recall, Multi-hop Evidence Completeness\\
Temporal & Temporal Validity Accuracy, Temporal Leakage Rate, Current/Historical Separation Accuracy, Comparison Separation Accuracy\\
Citation & Precision, Recall, F1, Invalid Citation Rate\\
Grounding & Numeric Grounding Accuracy, Unsupported Claim Rate, Claim Support Precision, Answer Evidence Completeness\\
Abstention & Precision, Recall, F1\\
Vận hành & P50/P95 latency, token usage, cost, parser time, indexing time\\
\bottomrule
\end{tabularx}
\end{table}

\section{Chống leakage và bảo đảm tái lập}

Final test set được đóng băng trước final run và không bao giờ dùng để tuning. Thay đổi gold label phải tạo version mới, không sửa file đã đóng băng. Mỗi run lưu run\_id, git commit, corpus version/hash, gold-set version/hash, suite và variant, run\_manifest\_hash, model IDs, prompt versions, parser versions, config snapshot, token usage, estimated cost và raw results path. Run và raw artifact bất biến/append-only; mọi query fail vẫn được giữ trong error analysis.

\section{Đe dọa đến tính hợp lệ}

Các đe dọa gồm corpus nhỏ và một reviewer chính có thể tạo annotation bias, LLM judge variance, provider model drift, query distribution chưa đại diện toàn bộ người dùng, và khả năng label gold sai. Biện pháp giảm thiểu là dùng metric xác định làm headline, pin model snapshot, lưu raw output, review độc lập label gold, chia data split rõ ràng và công khai limitation.

% ============================================================
\chapter{Thiết Kế Hệ Thống}
% ============================================================

\section{Nguyên tắc thiết kế}

Hệ thống tuân theo các nguyên tắc: parser-neutral IR, PostgreSQL là nguồn chân lý, Qdrant là index dẫn xuất dựng lại được, temporal-by-default, citation-by-ID, verified-or-abstain, structure-aware, chất lượng parser là mục tiêu hạng nhất, quan hệ tham chiếu được mô hình hóa tường minh, evidence completeness trước khi sinh câu trả lời, verification deterministic-first, failure-aware repair có giới hạn, controlled workflow (không autonomous agent), Langfuse ngoài đường tới hạn, RAGFlow chỉ là baseline bên ngoài, không dùng open-web search, reproducible experiments và local-first defense.

\section{Kiến trúc tổng quan}

\begin{figure}[H]
\centering
\resizebox{\textwidth}{!}{%
\begin{tikzpicture}[
  node distance=8mm and 10mm,
  box/.style={draw, rounded corners, minimum width=3.5cm, minimum height=0.9cm, align=center},
  arrow/.style={-{Latex[length=2mm]}, thick}
]
\node[box, fill=SoftBlue] (web) {Next.js UI};
\node[box, fill=SoftBlue, right=of web] (api) {FastAPI};
\node[box, fill=SoftGray, right=of api] (graph) {LangGraph Workflow};
\node[box, fill=SoftGray, right=of graph] (gen) {Generator + Verifier};

\node[box, fill=SoftBlue, below=1.3cm of api] (router) {Parser Router};
\node[box, fill=SoftGray, right=of router] (ir) {Canonical IR};
\node[box, fill=SoftGray, right=of ir] (lse) {Legal Extractor};
\node[box, fill=SoftBlue, right=of lse] (pg) {PostgreSQL};
\node[box, fill=SoftBlue, right=of pg] (qd) {Qdrant};

\draw[arrow] (web) -- (api);
\draw[arrow] (api) -- (graph);
\draw[arrow] (graph) -- (gen);
\draw[arrow] (api) -- (router);
\draw[arrow] (router) -- (ir);
\draw[arrow] (ir) -- (lse);
\draw[arrow] (lse) -- (pg);
\draw[arrow] (pg) -- (qd);
\draw[arrow] (qd) |- (graph);
\draw[arrow] (pg) |- (gen);
\end{tikzpicture}%
}
\caption{Kiến trúc logic tổng thể của hệ thống}
\label{fig:architecture}
\end{figure}

Hệ thống gồm hai pipeline tách biệt: offline ingestion và online query. Observability (Langfuse) chạy xuyên suốt nhưng không nằm trên đường tới hạn tính đúng đắn.

\section{Pipeline ingestion ngoại tuyến}

\begin{figure}[H]
\centering
\resizebox{0.95\textwidth}{!}{%
\begin{tikzpicture}[
  node distance=6mm,
  box/.style={draw, rounded corners, minimum width=4cm, minimum height=0.78cm, align=center, fill=SoftBlue},
  decision/.style={draw, diamond, aspect=2.2, align=center, fill=DraftYellow},
  arrow/.style={-{Latex[length=2mm]}, thick}
]
\node[box] (source) {PDF + manifest chính thức};
\node[box, below=of source] (queue) {Ingestion Queue (Redis + Dramatiq)};
\node[box, below=of queue] (router) {Parser Router (Docling | MinerU)};
\node[box, below=of router] (ir) {Canonical Document IR};
\node[box, below=of ir] (extract) {Legal Structure Extractor};
\node[box, below=of extract] (enrich) {Legal Context Enricher};
\node[box, below=of enrich] (resolve) {Reference + Temporal Resolver};
\node[decision, below=8mm of resolve] (route) {Quality Gates};
\node[box, below left=9mm and 13mm of route, fill=SoftGray] (review) {Needs review};
\node[box, below=9mm of route] (accept) {Accepted};
\node[box, below right=9mm and 13mm of route, fill=SoftGray] (drop) {Dropped};
\node[box, below=of review] (human) {Human Review};
\node[box, below=of human] (accepted2) {Accepted sau review};
\node[box, right=13mm of human, fill=SoftGray] (reject) {Rejected / audit};
\node[box, below=of accept] (store) {PostgreSQL -> Qdrant};
\draw[arrow] (source) -- (queue);
\draw[arrow] (queue) -- (router);
\draw[arrow] (router) -- (ir);
\draw[arrow] (ir) -- (extract);
\draw[arrow] (extract) -- (enrich);
\draw[arrow] (enrich) -- (resolve);
\draw[arrow] (resolve) -- (route);
\draw[arrow] (route) -- (review);
\draw[arrow] (route) -- (accept);
\draw[arrow] (route) -- (drop);
\draw[arrow] (review) -- (human);
\draw[arrow] (human) -- node[above, font=\small]{accept} (accepted2);
\draw[arrow] (human) -- node[above, font=\small]{reject} (reject);
\draw[arrow] (drop) -- (reject);
\draw[arrow] (accept) -- (store);
\draw[arrow] (accepted2) |- (store);
\end{tikzpicture}%
}
\caption{Pipeline xử lý và review tài liệu}
\label{fig:ingestion}
\end{figure}

\begin{sloppypar}
Pipeline worker thực hiện \code{parse -> normalize -> legal extract -> reference resolve -> temporal resolve -> quality gates -> review -> embed -> index}, mỗi bước là một actor Dramatiq ngắn, rời rạc và idempotent. Chỉ kết quả phân loại \code{accepted} mới được index; \code{needs\_review} phải qua quyết định reviewer; \code{dropped} chỉ được ghi audit và không bao giờ index. Qdrant chỉ nhận dữ liệu có \code{review\_status = ACCEPTED} đọc từ PostgreSQL.
\end{sloppypar}

\section{Parser Router và Canonical Document IR}

Parser Router quyết định parser theo đặc tính tài liệu và quality gate (FR-01). Docling là parser chính; MinerU là parser phụ và fallback/challenger. Quy tắc routing:

\begin{itemize}
  \item PDF searchable có text layer và layout chuẩn: Docling chạy trước;
  \item PDF scan hoặc layout lỗi: Docling chạy trước với OCR backend CPU, nếu quality gate fail thì chuyển MinerU;
  \item bảng phức tạp: so sánh đầu ra hai parser khi cần.
\end{itemize}

\begin{sloppypar}
Quality gate chia hai nhóm: nhóm A (parser-level) chạy sau khi chuẩn hóa IR với các chỉ số như provenance coverage, text extraction rate, table detection và layout coherence; nhóm B (structural) chạy sau Legal Structure Extractor với các chỉ số như point label detection, hierarchy completeness và short-Point retention. Nếu nhóm structural fail, kết quả structural hiện tại bị hủy và Router chạy lại toàn bộ pipeline từ parser thay thế, không trộn kết quả hai parser. Mọi quyết định routing và kết quả quality gate được ghi vào \code{ingestion\_runs\allowbreak.parser\_routing} và \code{DocumentElement\allowbreak.source\_parser}.
\end{sloppypar}

\begin{sloppypar}
Canonical Document IR là biểu diễn trung gian do dự án sở hữu, gồm \code{ParsedDocument} chứa \code{ParsedPage[]}, mỗi trang chứa \code{DocumentElement[]}. Mỗi element mang \code{element\_id}, \code{element\_type}, \code{text}, \code{page\_number}, \code{bbox}, \code{reading\_order}, \code{parent\_element\_id}, \code{table\_html}, \code{source\_parser}, \code{parser\_version}, \code{parser\_confidence} và \code{raw\_reference}. Tầng IR cô lập phân tích pháp lý khỏi định dạng đầu ra của Docling/MinerU.
\end{sloppypar}

\section{Legal Structure Extractor}

Legal Structure Extractor là parser pháp lý riêng của VNLRAG, chạy trên Canonical Document IR, nhận diện Chương, Mục, Điều, Khoản, Điểm, Phụ lục, bảng pháp lý, điều khoản chuyển tiếp, tiêu đề và đánh số văn bản pháp luật Việt Nam kèm biến thể do OCR.

Điểm nhấn quan trọng nhất là nhãn Điểm tiếng Việt: bắt buộc hỗ trợ a) b) c) d) đ) e) theo bảng chữ cái tiếng Việt 29 ký tự, không dùng giả định \code{[a-z]}. Ánh xạ sang provision ID giữ nguyên ký tự tiếng Việt: Điểm d) ứng \code{diem-d}, Điểm đ) ứng \code{diem-đ}, tránh va chạm ID. Khi OCR không phân biệt được \code{d)} và \code{đ)}, extractor dùng pattern ngữ cảnh; nếu không đủ ngữ cảnh thì gắn cờ ambiguity và định tuyến \code{needs\_review}, không tự chọn bừa.

Extractor cũng bảo đảm short-Point retention: một Điểm pháp lý ngắn nhưng hợp lệ vẫn là provision hợp lệ, không bị loại bỏ vì số token thấp. Quy tắc tạo provision ID:

\begin{center}
\code{\{loai-van-ban\}\allowbreak-\{so\}\allowbreak-\{nam\}\allowbreak\_\_dieu\allowbreak-\{n\}\allowbreak\_\_khoan\allowbreak-\{n\}\allowbreak\_\_diem\allowbreak-\{chu-cai\}}
\end{center}

\begin{sloppypar}
Ví dụ: \code{nd-168-2024\_\_dieu-7\_\_khoan-4\_\_diem-b}. Các node không thuộc cây Điều thường dùng dạng \code{\_\_phu-luc-\{n\}}, \code{\_\_phu-luc-\{n\}\_\_bang-\{m\}}, \code{\_\_dieu-\{n\}\_\_bang-\{m\}}, \code{\_\_dieu-\{n\}\_\_khoan-chuyen-tiep}, \code{\_\_chuyen-tiep-\{k\}}, \code{\_\_tieu-de-\{n\}}.
\end{sloppypar}

Legal Context Enricher bổ sung parent-context vào \code{retrieval\_text} (câu mở đầu của Khoản, tiêu đề Điều) để phục vụ retrieval; \code{source\_text} không bao giờ bị biến đổi. Trích dẫn vẫn trỏ tới provision thực tế. Ví dụ minh họa: \code{source\_text} là ``p) Dàn hàng ngang từ 03 xe trở lên''; \code{retrieval\_text} là ``Khoản 4. Phạt tiền từ ... đến ... đối với một trong các hành vi sau: p) Dàn hàng ngang từ 03 xe trở lên''. Các con số trong ví dụ là placeholder minh họa, không phải khẳng định nội dung văn bản thực tế.

\section{Mô hình dữ liệu}

\subsection{LegalDocument và DocumentVersion}

LegalDocument quản lý document ID, số hiệu, tiêu đề, loại văn bản, cơ quan ban hành, ngày ban hành, nguồn, file hash và trạng thái. DocumentVersion ghi nhận mỗi phiên bản nội dung với manifest bất biến, hash và khoảng hiệu lực, gắn với các quan hệ pháp lý.

\subsection{LegalProvision}

\begin{lstlisting}[language=Python,caption={Mô hình LegalProvision rút gọn}]
class LegalProvision(BaseModel):
    provision_id: str
    document_version_id: str
    node_kind: str                      # ARTICLE/CLAUSE/POINT/APPENDIX/TABLE/...
    chapter: str | None = None
    section: str | None = None
    article: str | None = None
    clause: str | None = None
    point: str | None = None
    heading: str | None = None
    source_text: str                    # bat bien, giu nguyen van ban goc
    retrieval_text: str                 # co the ke thua parent_context
    parent_context: str | None = None
    effective_from: date | None = None
    effective_to: date | None = None
    status: str
    page_number: int
    bbox: BoundingBox | None = None
    source_element_ids: list[str]
    content_hash: str
    version: int
    review_status: str                  # PENDING/ACCEPTED/REJECTED/DROPPED
\end{lstlisting}

\begin{sloppypar}
\code{source\_text} là nội dung pháp lý thuộc trực tiếp provision và giữ nguyên; \code{retrieval\_text} phục vụ retrieval và có thể kế thừa ngữ cảnh cha. \code{provision\_id} ổn định và tái tạo được; khi provision bị sửa đổi, ID giữ nguyên để trích dẫn ổn định, nội dung mới lưu dưới version mới. Unique key vật lý là cặp \code{(provision\_id, version)}. \code{status} mô tả vòng đời pháp lý của provision; \code{review\_status} mô tả trạng thái review corpus và là cổng chặn trong điều kiện hiệu lực; hai trường này độc lập nhau.
\end{sloppypar}

\subsection{ProvisionReference và DocumentRelation}

\begin{sloppypar}
ProvisionReference ghi quan hệ cấp provision với bốn loại: \code{PARENT\_OF}, \code{REFERS\_TO}, \code{SIBLING\_OF}, \code{PENALTY\_COMPANION}. Trong đó \code{PENALTY\_COMPANION} gắn quy định xử phạt với quy định đi kèm như trừ điểm giấy phép lái xe hoặc tước quyền sử dụng giấy phép lái xe. DocumentRelation ghi quan hệ cấp văn bản với sáu loại: \code{AMENDS}, \code{REPEALS}, \code{SUPERSEDES}, \code{CORRECTS}, \code{GUIDES}, \code{RELATED\_TO}. Quan hệ được version-bound (\code{source\_provision\_version\_id}, \code{target\_provision\_version\_id}) và lưu trong bảng PostgreSQL, xử lý bằng application logic; không dùng Neo4j hoặc knowledge graph. Reference không giải quyết được ghi \code{UNRESOLVED}/\code{PENDING\_REVIEW} và định tuyến review, không suy đoán quan hệ.
\end{sloppypar}

\subsection{LegalEffectEvent}

\begin{sloppypar}
LegalEffectEvent ghi lại sự kiện pháp lý ảnh hưởng tới hiệu lực của văn bản/provision: \code{EFFECTIVE}, \code{AMENDED}, \code{PARTIAL\_AMENDED}, \code{SUPERSEDED}, \code{REPEALED}, \code{CORRECTED}, \code{EXPIRED}. \code{affected\_provision\_versions} liệt kê structured các provision/version chịu ảnh hưởng. Sự kiện này phục vụ Temporal and Amendment Resolver và câu hỏi so sánh lịch sử.
\end{sloppypar}

\section{PostgreSQL schema}

PostgreSQL là nguồn chân lý, lưu toàn bộ metadata, phiên bản, quan hệ, review, audit, query trace và feedback \cite{postgresql2026}.

\begin{table}[H]
\caption{Các bảng dữ liệu chính}
\begin{tabularx}{\textwidth}{p{4.5cm}X}
\toprule
\textbf{Bảng} & \textbf{Vai trò}\\
\midrule
legal\_sources & Nguồn văn bản chính thống và mức ưu tiên đối chiếu\\
legal\_documents & Document và metadata\\
document\_versions & Phiên bản văn bản, manifest, hash và khoảng hiệu lực\\
legal\_provisions & Provision version đầy đủ: hierarchy, text, interval, review\\
provision\_versions & Registry lineage, \code{superseded\_by\_version}\\
provision\_references & Quan hệ cấp provision, version-bound\\
document\_relations & Quan hệ cấp văn bản (AMENDS, SUPERSEDES, REPEALS, ...)\\
legal\_effect\_events & Sự kiện pháp lý ảnh hưởng hiệu lực\\
parsed\_documents / document\_elements & Canonical Document IR đã chuẩn hóa\\
ingestion\_runs & Trạng thái job và parser routing\\
ingestion\_artifacts & Object key trong MinIO và metadata\\
review\_items & Bằng chứng và quyết định reviewer\\
query\_traces / query\_feedback & Lịch sử query, verification summary và feedback\\
evaluation\_datasets / evaluation\_runs / evaluation\_results & Gold set, run metadata và raw result\\
corpus\_qa\_reports & Báo cáo chất lượng corpus\\
\bottomrule
\end{tabularx}
\end{table}

\begin{sloppypar}
Các ràng buộc quan trọng: UNIQUE \code{(provision\_id, version)}; CHECK interval \code{effective\_to IS NULL OR effective\_to > effective\_from}; CHECK review-required \code{review\_status <> 'ACCEPTED' OR effective\_from IS NOT NULL}; exclusion constraint daterange ngăn hai version ACCEPTED chồng lấn trong cùng provision. Alembic quản lý migration; Qdrant dựng lại được từ \code{SELECT * FROM legal\_provisions WHERE review\_status = 'ACCEPTED'}.
\end{sloppypar}

\section{Qdrant collection}

Mỗi Qdrant point đại diện cho một provision version ACCEPTED. Collection có named dense vector \code{dense} (768 chiều khởi điểm, Cosine) và sparse vector \code{sparse} (BM25). Payload chứa provision ID, provision version, document ID/version, number, type, hierarchy fields, vehicle types, interval, review status, parser version, legal parser version, sparse encoder version, content hash và relation metadata giới hạn (\code{REFERS\_TO}/\code{PENALTY\_COMPANION} đã RESOLVED và ACCEPTED). Qdrant hỗ trợ hybrid query, RRF và payload filter trực tiếp \cite{qdrant2026hybrid}.

Collection được version hóa, ví dụ \code{legal\_provisions\_v1}. Alias \code{legal\_provisions\_active} trỏ đến version đang phục vụ. Khi đổi embedding, sparse encoder, payload schema hoặc chunking production, hệ thống tạo collection mới, chạy regression rồi switch alias; không bao giờ rebuild in place và không trộn hai embedding space.

\section{Retrieval đa tầng}

\subsection{Query Understanding và temporal resolution}

Query Understanding ưu tiên deterministic parsing bằng regex cho ngày, năm, số hiệu văn bản, Điều/Khoản/Điểm và loại phương tiện; LLM structured extraction chỉ xử lý phần còn lại. Output là QueryPlan gồm intent, effective date, comparison dates, vehicle type, document number, article/clause/point, legal entities, normalized query, required evidence (evidence plan) và missing query information. Các intent gồm CURRENT, HISTORICAL, COMPARISON, SOURCE\_SEARCH và OUT\_OF\_SCOPE.

Chính sách ngày: không có ngày thì dùng ngày request; có ngày cụ thể thì dùng ngày đó; chỉ có năm và không có sự kiện đổi hiệu lực trong năm thì áp dụng canonical date (ví dụ 01/07 của năm) và bắt buộc hiển thị ngày đã áp dụng; có sự kiện đổi hiệu lực trong năm thì yêu cầu ngày cụ thể hoặc ABSTAIN với \code{MISSING\_QUERY\_DATE}. Hệ thống không dùng văn bản hiện hành làm mặc định cho câu hỏi lịch sử.

\subsection{Query Expansion}

Query expansion luôn giữ câu hỏi gốc của người dùng. Các variant gồm original, normalized (chuẩn hóa thuật ngữ pháp lý như ``phạt bao nhiêu'' sang ``mức xử phạt''), multi-query rewrite (giới hạn, không đệ quy) và conditional HyDE. HyDE chỉ bật khi câu ngắn, khẩu ngữ, ngữ nghĩa yếu hoặc bằng chứng chưa đủ, không bật luôn và không dùng cho exact lookup hay sparse.

\subsection{Ba kênh retrieval, RRF và reranking}

Ba kênh chạy song song:

\begin{enumerate}
  \item \begin{sloppypar} exact legal lookup theo payload filter và SQL cho định danh như \code{168/2024/NĐ-CP}, \code{Điều 7}, \code{Khoản 4}, \code{Điểm a}; \end{sloppypar}
  \item dense semantic với named vector \code{dense};
  \item sparse lexical BM25 với named vector \code{sparse}.
\end{enumerate}

Dense và sparse được hợp nhất bằng RRF (k và weights cấu hình được), sau đó post-filter temporal invariant và dedup theo \code{provision\_id}; các candidate exact lookup được bảo toàn và ghép vào tập đã fusion theo policy rõ ràng (exact match được ưu tiên giữ lại, loại trùng lặp theo \code{provision\_id}) trước khi rerank. Reranking dùng Jina Reranker v3 như stage chuẩn của pipeline, nhưng không khẳng định reranker cải thiện chất lượng trước khi có kết quả R6 của Suite C.

\subsection{Legal context expansion}

Mở rộng ngữ cảnh chỉ quanh các seed provision mạnh theo parent, sibling, direct cross-reference và penalty companion, cùng các quy định liên quan sửa đổi/thay thế. Mỗi provision mở rộng ghi lý do vào context:

\begin{lstlisting}[language=json,caption={Metadata của provision được mở rộng}]
{
  "provision_id": "nd-168-2024__dieu-7__khoan-4__diem-b",
  "added_by": "CROSS_REFERENCE",
  "source_id": "nd-168-2024__dieu-33__khoan-2",
  "depth": 1
}
\end{lstlisting}

Mở rộng theo quan hệ phải áp dụng temporal filter và review filter theo ngày query; quan hệ PENDING\_REVIEW/UNRESOLVED không được dùng để mở rộng tự động; depth có giới hạn để tránh mở rộng đồ thị vô hạn.

\section{Evidence Completeness Gate}

Trước khi sinh câu trả lời, Evidence Completeness Gate kiểm tra mọi loại bằng chứng trong evidence plan đã được thu thập trong context. Nếu trạng thái là \code{INCOMPLETE} (ví dụ câu hỏi mức phạt + điểm trừ mà chỉ tìm được mức phạt), hệ thống chạy targeted retrieval theo \code{evidence\_gaps} hoặc mở rộng ngữ cảnh theo quan hệ, rồi kiểm tra lại trước khi gọi generator. Mỗi lượt quay lại tính vào \code{repair\_attempts} hữu hạn; hết giới hạn thì ABSTAIN. Hệ thống không âm thầm trả lời một nửa dễ của câu hỏi đa bằng chứng.

\section{Workflow hỏi đáp}

LangGraph phân biệt workflow và agent: workflow có code path định trước, còn agent tự quyết định process/tool usage \cite{langgraph2026workflows}. Đề tài dùng workflow có kiểm soát với các nhánh xác định trước, không phải autonomous multi-agent.

\begin{figure}[H]
\centering
\resizebox{\textwidth}{!}{%
\begin{tikzpicture}[
  node distance=7mm and 10mm,
  box/.style={draw, rounded corners, minimum width=3.1cm, minimum height=0.85cm, align=center, fill=SoftBlue},
  decision/.style={draw, diamond, aspect=2.1, align=center, fill=DraftYellow},
  arrow/.style={-{Latex[length=2mm]}, thick}
]
\node[box] (analyze) {Analyze query};
\node[box, right=of analyze] (temporal) {Resolve temporal};
\node[box, right=of temporal] (expand) {Expand query};
\node[box, right=of expand] (retrieve) {Parallel recall};
\node[box, right=of retrieve] (fuse) {Fuse + rerank};
\node[box, right=of fuse] (lce) {Context expansion};
\node[decision, right=of lce] (eg) {Evidence gate};
\node[box, below=14mm of eg, fill=SoftGray] (target) {Targeted retrieval};
\node[box, right=14mm of eg] (ctx) {Build context};
\node[box, right=of ctx] (gen) {Structured generation};
\node[decision, right=of gen] (verify) {Verify L1-L6};
\node[box, above right=7mm and 10mm of verify] (final) {Verified answer};
\node[box, below right=7mm and 10mm of verify, fill=SoftGray] (abstain) {Abstain};
\draw[arrow] (analyze) -- (temporal);
\draw[arrow] (temporal) -- (expand);
\draw[arrow] (expand) -- (retrieve);
\draw[arrow] (retrieve) -- (fuse);
\draw[arrow] (fuse) -- (lce);
\draw[arrow] (lce) -- (eg);
\draw[arrow] (eg) -- node[above, font=\small]{complete} (ctx);
\draw[arrow] (eg) -- node[right, font=\small]{incomplete} (target);
\draw[arrow] (target) |- (eg);
\draw[arrow] (ctx) -- (gen);
\draw[arrow] (gen) -- (verify);
\draw[arrow] (verify) -- (final);
\draw[arrow] (verify) -- (abstain);
\end{tikzpicture}%
}
\caption{Controlled workflow cho online query}
\label{fig:query-workflow}
\end{figure}

\begin{sloppypar}
Đồ thị đề xuất gồm: \code{analyze\_query}, \code{resolve\_temporal}, \code{expand\_query}, \code{retrieve\_parallel}, \code{fuse}, \code{rerank}, \code{expand\_legal\_context}, \code{check\_evidence}, \code{targeted\_retrieval}, \code{build\_context}, \code{generate}, \code{verify}, \code{finalize}, \code{repair} và \code{abstain}. Các cạnh có điều kiện: \code{check\_evidence} phân nhánh \code{complete} sang \code{build\_context} rồi \code{generate}, hoặc \code{incomplete} sang \code{targeted\_retrieval} rồi quay lại \code{check\_evidence}; \code{verify} phân nhánh \code{valid} sang \code{finalize}, \code{repairable} sang đường repair theo loại lỗi, còn lại sang \code{abstain}. Mọi đường quay lại đều tăng \code{repair\_attempts} trong state; khi \code{repair\_attempts >= max\_repair\_attempts} (config, khởi điểm 3) thì ABSTAIN. Không tồn tại vòng lặp không tăng counter.
\end{sloppypar}

\section{Structured generation}

Generator (Gemini 3.5 Flash) sinh câu trả lời theo schema cấp claim:

\begin{lstlisting}[language=json,caption={Structured output của generator}]
{
  "answer_summary": "...",
  "claims": [
    {
      "claim": "Người điều khiển xe mô tô vượt đèn đỏ bị xử phạt tiền.",
      "claim_type": "MONETARY_PENALTY",
      "provision_ids": [
        "nd-168-2024__dieu-7__khoan-4__diem-b"
      ],
      "numbers": ["4.000.000", "6.000.000"]
    }
  ],
  "missing_information": [],
  "should_abstain": false
}
\end{lstlisting}

\begin{sloppypar}
Schema dùng Pydantic với \code{extra="forbid"}, \code{ClaimType} là enum đóng, \code{claims} bắt buộc tối thiểu một claim và mỗi \code{provision\_ids} có tối thiểu một ID khi \code{should\_abstain=false}. Các con số trong ví dụ mang tính minh họa cấu trúc dữ liệu, không phải khẳng định giá trị thực tế của văn bản. Model không được viết citation text; phần trích dẫn hiển thị được dựng từ metadata đã lưu, không phải chuỗi LLM gõ tự do. Khi draft có \code{should\_abstain=true}, workflow không chuyển sang \code{finalize} mà route sang abstention với reason code ánh xạ và bảo toàn \code{missing\_information}.
\end{sloppypar}

\section{Verification sáu tầng}

Verifier chạy theo thứ tự sáu tầng, deterministic-first; LLM judge độc lập chỉ được gọi ở L5 cho trường hợp ngữ nghĩa:

\begin{table}[H]
\caption{Sáu tầng verification}
\begin{tabularx}{\textwidth}{p{1.2cm}p{4.2cm}Y}
\toprule
\textbf{Tầng} & \textbf{Bộ kiểm chứng} & \textbf{Nội dung}\\
\midrule
L1 & Schema verifier & Output tuân thủ Pydantic schema cấp claim; không unknown field; summary ánh xạ atomic sang claim\\
L2 & Citation ID verifier & \code{provision\_id} tồn tại; thuộc context whitelist hoặc mở rộng hợp lệ; \code{review\_status = ACCEPTED}; metadata citation có thẩm quyền\\
L3 & Temporal verifier & Provision trích dẫn hợp lệ tại ngày được hỏi\\
L4 & Numeric grounding verifier & Mức phạt, số điểm trừ, ngày, tuổi, thời hạn, số lượng khớp giá trị bằng chứng đã chuẩn hóa\\
L5 & Claim support verifier & Deterministic trước; LLM judge độc lập chỉ cho trường hợp ngữ nghĩa, fail-closed\\
L6 & Evidence completeness verifier & Mọi loại bằng chứng trong evidence plan được claims cuối bao phủ\\
\bottomrule
\end{tabularx}
\end{table}

Aggregate validity chỉ đúng khi cả sáu tầng đều pass, không có ``pass mềm''. Bất biến API: Returned Invalid Citation Rate = 0. Nếu hệ thống không thể bảo đảm invariant này, request phải abstain hoặc trả lỗi.

Failure-aware repair xử lý lỗi theo loại cụ thể, không chỉ regenerate: thiếu bằng chứng chạy targeted retrieval rồi dựng lại context và regenerate; claim không được hỗ trợ regenerate từ bằng chứng hiện có hoặc targeted retrieval nếu thiếu; schema sai regenerate structured output; xung đột thời gian truy xuất phiên bản thời gian đúng. Mọi nhánh repair cùng tính vào \code{MAX\_REPAIR\_ATTEMPTS}; hết giới hạn thì ABSTAIN.

\section{Abstention}

Các reason code chuẩn gồm:
\begin{itemize}
  \item \code{OUT\_OF\_SCOPE}, \code{MISSING\_QUERY\_DATE};
  \item \code{INSUFFICIENT\_EVIDENCE}, \code{NO\_VALID\_PROVISION};
  \item \code{TEMPORAL\_CONFLICT}, \code{CITATION\_VERIFICATION\_FAILED};
  \item \code{CORPUS\_NOT\_COVERED}.
\end{itemize}

Response abstention không chứa answer pháp lý, citations rỗng, có disclaimer, có trace ID, có message rõ ràng và thông tin cần bổ sung. Hệ thống không hiển thị confidence phần trăm nếu chưa được calibration.

\section{Langfuse observability}

\begin{sloppypar}
Langfuse phục vụ tracing, token usage, cost, latency, prompt management và experiments. Trace \code{legal\_query} gồm các span \code{analyze\_query}, \code{normalize\_query}, \code{rewrite\_query}, \code{hyde}, \code{exact\_lookup}, \code{dense\_retrieval}, \code{sparse\_retrieval}, \code{rrf\_fusion}, \code{reranker}, \code{reference\_expansion}, \code{evidence\_check}, \code{generate}, \code{citation\_verify}, \code{numeric\_verify}, \code{claim\_verify}.
\end{sloppypar}

Langfuse \textbf{không nằm trên đường tới hạn tính đúng đắn}: trace được ingest bất đồng bộ, toàn bộ callback non-mutating; nếu Langfuse không khả dụng, query vẫn hoạt động bình thường. Observability giải thích hệ thống chạy như thế nào và chi phí ra sao, nhưng không thay thế logic correctness - tính đúng đắn do retrieval, evidence gate và verifier L1-L6 bảo đảm. Prompt được quản lý theo version và label production/dev; version prompt được ghi vào từng evaluation run và \code{QueryTrace.config\_snapshot}.

\section{Background ingestion và object storage}

\begin{sloppypar}
Ingestion chạy hoàn toàn qua hàng đợi: \apipath{POST /api/v1/documents} trả \code{202 Accepted} kèm \code{ingestion\_job\_id}; không parse PDF đồng bộ trong request handler. Redis làm broker cho Dramatiq và cache; job state nằm trong PostgreSQL (\code{ingestion\_runs}) là source of truth. Mỗi bước là một actor idempotent ngắn (\code{parse}, \code{normalize}, \code{extract}, \code{resolve\_refs}, \code{resolve\_temporal}, \code{quality\_gate}, \code{embed}, \code{index}) có time limit cấu hình riêng; retry chỉ áp dụng cho lỗi transient. Worker fail được phục hồi bằng retry, dead-letter queue và script reconcile giữa job state và index state.
\end{sloppypar}

\begin{sloppypar}
MinIO là object storage S3-compatible, lưu PDF nguồn, đầu ra parser, ảnh trang và artifact ingestion/review/evaluation trong các bucket riêng (\code{source-pdfs}, \code{parser-outputs}, \code{page-images}, \code{ingestion-artifacts}, \code{review-artifacts}, \code{evaluation-artifacts}). PostgreSQL lưu object key và metadata; backup MinIO dùng replication hoặc \code{mc mirror} sang nơi lưu trữ độc lập, không dùng tiering/ILM làm backup.
\end{sloppypar}

\section{API và giao diện}

FastAPI cung cấp chat, search, provision, documents, jobs, reviews, feedback, health, evaluations và corpus-qa endpoints \cite{fastapi2026}. Mọi response nghiệp vụ chứa \code{trace\_id}; abstention dùng HTTP 200 kèm trạng thái \code{ABSTAINED}; lỗi kỹ thuật dùng 4xx/5xx.

\begin{table}[H]
\caption{API chính}
\small
\begin{tabularx}{\textwidth}{p{1.6cm}p{5.2cm}Y}
\toprule
\textbf{Method} & \textbf{Path} & \textbf{Vai trò}\\
\midrule
POST & \apipath{/api/v1/chat} & Verified answer hoặc abstention\\
POST & \apipath{/api/v1/search} & Search provision, không gọi generator\\
POST & \apipath{/api/v1/documents} & Upload PDF + manifest, trả 202 + job ID\\
GET & \apipath{/api/v1/jobs/\{job\_id\}} & Trạng thái job ingestion\\
\makecell{GET\\POST} & \apipath{/api/v1/reviews} & Danh sách và quyết định review\\
POST & \apipath{/api/v1/feedback} & Feedback người dùng cuối\\
POST & \apipath{/api/v1/evaluations} & Chạy evaluation run\\
GET & \apipath{/api/v1/corpus-qa/report} & Báo cáo chất lượng corpus\\
GET & \apipath{/api/v1/health/ready} & Readiness PostgreSQL/Qdrant/Redis/MinIO\\
\bottomrule
\end{tabularx}
\end{table}

Frontend Next.js có chat page, search page, citation card dựng từ metadata, passage viewer mở snippet nguồn kèm trang, date selector hiển thị ngày áp dụng, abstention panel và feedback widget. Giao diện không stream raw token vì draft có thể chưa pass verifier; SSE chỉ stream progress event như query analyzed, retrieval completed, generation started và verification completed.

\section{Bảo mật}

Admin endpoint dùng Bearer token trong P0. Upload được kiểm tra MIME, magic bytes, size, filename và SHA-256 duplicate; filename nội bộ do hệ thống sinh để chặn path traversal. PDF text được coi là data, không phải instruction, nhằm giảm prompt injection; nội dung ``ignore previous instructions'' không được làm thay đổi hành vi và claim không được hỗ trợ bị verifier từ chối. API key không được log hoặc đưa vào image. Rate limiting cấu hình theo deployment, không hardcode theo free-tier quota. Hệ thống chỉ mô tả biện pháp giảm thiểu dữ liệu cá nhân, không tuyên bố tuân thủ pháp lý tuyệt đối khi chưa có compliance review \cite{nghidinh13_2023}.

% ============================================================
\chapter{Cài Đặt và Triển Khai}
% ============================================================

\section{Tech stack}

\begin{table}[H]
\caption{Công nghệ triển khai}
\begin{tabularx}{\textwidth}{p{4cm}X}
\toprule
\textbf{Thành phần} & \textbf{Công nghệ}\\
\midrule
Ngôn ngữ/backend & Python 3.11, FastAPI, Pydantic v2\\
Workflow & LangGraph 1.x controlled workflow\\
Parser chính & Docling 2.x\\
Parser phụ/fallback & MinerU 3.4.x (Parser Router)\\
Metadata/versioning & PostgreSQL 18, SQLAlchemy 2, Alembic\\
Retrieval & Qdrant v1.19 dense + sparse + RRF + payload filter\\
Dense embedding & Ứng viên: Gemini Embedding 2, Jina Embeddings v5 text-nano/small\\
Reranker & Jina Reranker v3 (stage chuẩn)\\
Generator & Gemini 3.5 Flash\\
Judge độc lập & GPT-5.4 mini snapshot \cite{openai2026gpt54mini}\\
Background jobs & Redis 8 + Dramatiq 2.x\\
Object storage & MinIO (S3-compatible)\\
Observability & Langfuse Cloud\\
Evaluation & Deterministic metrics + Ragas v0.4.x\\
Frontend & Next.js, TypeScript, shadcn/ui\\
Deployment & Docker Compose, GitHub Actions\\
Package manager & uv\\
\bottomrule
\end{tabularx}
\end{table}

Model ID nằm trong cấu hình, không hardcode trong domain logic. Embedding production chưa được chốt vĩnh viễn; quyết định dựa trên bằng chứng Suite B. Không dùng pgvector trong thiết kế này: vector retrieval nằm trong Qdrant, PostgreSQL giữ metadata và quan hệ pháp lý.

\section{Cấu trúc source code}

Source được chia thành API, domain, ingestion, retrieval, query, workflow, generation, verification, persistence, evaluation, feedback và observability. Domain không import FastAPI, SQLAlchemy, Qdrant client hoặc provider SDK để giữ khả năng kiểm thử độc lập.

\begin{lstlisting}[caption={Cấu trúc backend rút gọn}]
backend/app/
  api/
  domain/
  ingestion/
    parser_router.py
    adapters/           # docling_adapter, mineru_adapter
    document_ir.py      # ParsedDocument, ParsedPage, DocumentElement
    structure_extractor.py
    context_enricher.py
    reference_resolver.py
    temporal_resolver.py
    quality_gates.py
    actors/             # parse, extract, resolve_refs, quality_gate, embed, index
    indexer.py
  retrieval/            # embedding, sparse, qdrant_store, hybrid, reranker
  query/                # query_understanding, expansion, hyde, evidence_gate
  workflow/             # graph.py, state.py, nodes/
  generation/           # provider, gemini, prompts, schemas
  verification/         # l1_schema ... l6_evidence
  persistence/          # database, repositories, models
  evaluation/           # runner, suites, deterministic_metrics, ragas_metrics
  feedback/
  observability/        # logging, tracing (Langfuse)
\end{lstlisting}

\section{Cài đặt ingestion}

Ingestion service nhận PDF và manifest, tính SHA-256, kiểm tra duplicate, lưu source file lên MinIO, enqueue \code{parse\_actor} và trả \code{202 Accepted}. Worker chạy pipeline parse, normalize, legal extract, reference resolve, temporal resolve, quality gates; kết quả \code{accepted} được lưu vào PostgreSQL rồi embed và upsert vào Qdrant. Qdrant upsert chạy sau PostgreSQL transaction commit; nếu upsert fail, job giữ trạng thái \code{INDEXING} và được retry bằng background/CLI reconcile, không rollback dữ liệu PostgreSQL. Actor được thiết kế idempotent: đọc state job từ PostgreSQL, bỏ qua bước đã hoàn thành.

\section{Cài đặt retrieval}

Exact lookup chuyển định danh như ``Nghị định 168, Điều 7'' thành payload filter. Dense và sparse prefetch được thực hiện song song trong Qdrant Query API; RRF fusion trả candidate; backend áp dụng temporal invariant lần cuối để correctness không phụ thuộc một filter implementation cụ thể. Config retrieval được version hóa:

\begin{lstlisting}[caption={Ví dụ retrieval config}]
retrieval:
  exact_lookup:
    enabled: true
    filter_fields: [document_number, article, clause, point]
  dense_prefetch: 30
  sparse_prefetch: 30
  rrf:
    k: 60
    weights: {dense: 1.0, sparse: 1.0}
  fusion_limit: 20
  final_top_k: 8
  temporal_filter: true
  dedup_key: provision_id
  reranker: jina-reranker-v3
\end{lstlisting}

Các con số là khởi điểm, chỉ được khóa sau ablation Suite C.

\section{Cài đặt verifier}

\begin{sloppypar}
Verifier không parse citation regex. Nó nhận Pydantic object và context map. Mỗi tầng L1-L6 là một module riêng trả \code{LayerResult} gồm \code{passed}, \code{issues}, \code{checked\_claim\_indexes}, \code{checked\_provision\_version\_ids} và \code{repair\_action}. L4 chuẩn hóa số tiền, điểm, ngày trước khi so sánh với giá trị bằng chứng. L5 chạy deterministic trước, LLM judge độc lập (GPT-5.4 mini snapshot) chỉ cho trường hợp ngữ nghĩa với hành vi fail-closed. Returned invalid citation là invariant bằng 0; nếu hệ thống không thể bảo đảm invariant này, request phải abstain hoặc error.
\end{sloppypar}

\section{Giao diện người dùng}

\placeholderfigure{Màn hình chat với bộ chọn ngày, loại phương tiện và citation panel}
\placeholderfigure{Màn hình so sánh quy định giữa hai thời điểm}
\placeholderfigure{Source drawer hiển thị passage, trang và hierarchy}

Citation card được dựng từ database metadata đã verify, không phải chuỗi LLM. Passage viewer mở snippet nguồn kèm trang dựa trên page number và provenance. Ngày áp dụng được hiển thị rõ trên mỗi answer. Giao diện không stream raw token vì draft có thể chưa pass verifier.

\section{Docker Compose}

\begin{sloppypar}
Bảy service chính: frontend (Next.js), backend (FastAPI), worker (Dramatiq), PostgreSQL, Qdrant, Redis và MinIO. Mọi service nằm trong một Docker network nội bộ; chỉ hai port public local là \code{127.0.0.1:3000} (frontend) và \code{127.0.0.1:8000} (backend). One-shot service \code{migrate} chạy \code{alembic upgrade head}, tạo collection Qdrant và bucket MinIO trước khi backend/worker/frontend khởi động. RAGFlow chạy trong môi trường benchmark riêng, không nằm trong compose production.
\end{sloppypar}

\begin{figure}[H]
\centering
\resizebox{0.95\textwidth}{!}{%
\begin{tikzpicture}[
  node distance=12mm and 14mm,
  box/.style={draw, rounded corners, minimum width=3.5cm, minimum height=1cm, align=center, fill=SoftBlue},
  arrow/.style={-{Latex[length=2mm]}, thick}
]
\node[box] (browser) {Browser};
\node[box, right=of browser] (front) {Next.js};
\node[box, right=of front] (back) {FastAPI};
\node[box, right=of back] (worker) {Dramatiq Worker};
\node[box, below left=12mm and 4mm of back] (pg) {PostgreSQL};
\node[box, below right=12mm and 4mm of back] (qd) {Qdrant};
\node[box, below=12mm of back] (redis) {Redis};
\node[box, below right=12mm and 4mm of qd] (minio) {MinIO};
\draw[arrow] (browser) -- (front);
\draw[arrow] (front) -- (back);
\draw[arrow] (back) -- (pg);
\draw[arrow] (back) -- (qd);
\draw[arrow] (back) -- (redis);
\draw[arrow] (back) -- (minio);
\draw[arrow] (worker) -- (redis);
\draw[arrow] (worker) -- (pg);
\draw[arrow] (worker) -- (qd);
\draw[arrow] (worker) -- (minio);
\end{tikzpicture}%
}
\caption{Kiến trúc triển khai local với bảy service}
\end{figure}

Các provider bên ngoài tùy chọn gồm Langfuse Cloud, Gemini API, OpenAI API (judge) và Jina API (embedding/reranker). Trong release, image bên thứ ba được pin exact tag và thay bằng digest trước code freeze.

\section{Migration, backup và restore}

\begin{sloppypar}
Alembic quản lý schema; migration chỉ chạy bởi one-shot service \code{migrate}, application process không bao giờ tự chạy \code{alembic upgrade head}. Release bundle gồm PostgreSQL dump, Qdrant snapshot (copy sang nơi lưu trữ độc lập), MinIO mirror, gold set, evaluation results, release manifest, git tag và checksum. Vì PostgreSQL là source of truth, Qdrant có thể rebuild từ \code{SELECT * FROM legal\_provisions WHERE review\_status = 'ACCEPTED'} bằng alias switch; snapshot chỉ rút ngắn RTO, không phải nguồn chân lý. Restore được kiểm chứng bằng script \code{verify\_release.py} so sánh số point, hash corpus và migration revision.
\end{sloppypar}

\section{CI/CD và quality gate}

PR pipeline chạy lint (ruff), type check (mypy), unit test với coverage core tối thiểu 80\%, integration test trên PostgreSQL/Qdrant/Redis/MinIO thật (gồm \code{alembic upgrade head} từ database rỗng), frontend build và Docker build. Full LLM evaluation chỉ chạy manual ở feature freeze và release candidate để kiểm soát chi phí và nondeterminism. Bất kỳ PR nào ảnh hưởng retrieval phải chạy retrieval regression subset, temporal regression, citation invariant và gold-set integrity.

\section{RAGFlow baseline (external)}

RAGFlow không phải nền tảng chính mà chỉ là baseline so sánh bên ngoài (FR-31). Bốn variant bắt buộc: RAGFlow default, RAGFlow + Docling, RAGFlow + MinerU và pipeline VNLRAG custom legal-aware, chạy trên cùng corpus và cùng bộ câu hỏi evaluation trong môi trường benchmark riêng. Chỉ số so sánh tối thiểu gồm Recall@10, citation correctness, temporal leakage và evidence completeness. RAGFlow yêu cầu tối thiểu 4 CPU, 16 GB RAM và 50 GB disk, không chạy cùng lúc với ingestion/demo/eval nặng trên cùng máy. Kết quả baseline không bao giờ được báo cáo như kết quả thực nghiệm của VNLRAG.

\section{Ràng buộc phần cứng}

Máy mục tiêu dùng Intel Core i5-1035G1, 19 GB RAM và NVIDIA MX330 2 GB VRAM. Online stack chạy CPU và gọi model qua API; GPU local không dùng được cho LLM, embedding hay reranker. Docling chạy CPU trong worker với \code{MAX\_INGESTION\_WORKERS=1}; MinerU chỉ dùng pipeline backend CPU, không bao giờ chạy VLM/hybrid local (cần GPU tối thiểu 8 GB VRAM). Nếu RAM thực tế vượt budget, MinerU chuyển sang remote \code{*-http-client} hoặc host tách biệt; quyết định được ghi vào tài liệu vận hành.

\section{Kế hoạch thời gian}

\begin{table}[H]
\caption{Mốc triển khai chính}
\begin{tabularx}{\textwidth}{p{3.2cm}p{3.4cm}X}
\toprule
\textbf{Thời gian} & \textbf{Mốc} & \textbf{Deliverable}\\
\midrule
20/07-26/07 & Chốt thiết kế và corpus & ADR, manifest, schema, Parser Router, Canonical Document IR\\
27/07-02/08 & Parser foundation & Suite A (P1-P3), adapter Docling/MinerU, quality gates\\
03/08-09/08 & Legal extraction và data platform & Legal Structure Extractor, corpus QA, PostgreSQL, Qdrant, MinIO, ingestion queue\\
10/08-16/08 & Relation graph và temporal & Reference Resolver, Temporal Resolver, parent-context enrichment\\
17/08-23/08 & Retrieval và query workflow & Suite B và C (R1-R10), query expansion, reranker, evidence gate, LangGraph skeleton\\
24/08-30/08 & Verification, Langfuse, frontend & Verifier L1-L6, LangGraph full flow, trace, chat UI, feedback\\
31/08-06/09 & Evaluation và stabilization & Gold set 200 câu, Suite D (G1-G7), RAGFlow baseline, feature freeze\\
07/09-12/09 & Finalization & Code freeze, final evaluation, report, release candidate\\
13/09-14/09 & Rehearsal/defense & Demo local và bảo vệ\\
\bottomrule
\end{tabularx}
\end{table}

Mốc kiểm soát cứng: feature freeze 06/09/2026, code freeze 10/09/2026, release candidate 12/09/2026. Không thêm tính năng mới trong hai ngày cuối; không thu hẹp phạm vi vì lịch trình.

% ============================================================
\chapter{Kiểm Thử và Đánh Giá}
% ============================================================

\section{Chiến lược kiểm thử}

Test pyramid gồm unit, integration, regression và E2E. Unit test tập trung manifest, parser, structure extractor, stable ID, temporal interval, verifier L1-L6, metric và abstention; coverage core deterministic tối thiểu 80\%. Integration test bắt buộc chạy trên PostgreSQL, Qdrant, Redis và MinIO thật (không dùng SQLite để kết luận integration pass). Regression test ngăn chất lượng retrieval/citation/temporal giảm. E2E kiểm tra current, historical, comparison, abstention, multi-evidence, numeric-grounding-fail và scan PDF qua Parser Router.

\section{Test extraction và parser}

Golden fixture bao phủ Điều không có Khoản, Điều có nhiều Khoản, Khoản có nhiều Điểm, Điểm kéo dài nhiều paragraph, Chương và Mục, heading xuống dòng, OCR spacing bất thường, số thứ tự giả trong bảng, Phụ lục, nhãn \code{d)} và \code{đ)} không bị lẫn, Điểm ngắn hợp lệ, header/footer lặp, điều khoản chuyển tiếp và d/đ ambiguity. Invariant bắt buộc: \code{source\_text} không rỗng, article order deterministic, không có hai provision trùng stable ID, d/đ ambiguity không tạo provision ID collision. Hierarchy extraction được chấm bằng tuple document/article/clause/point; fixture stable-ID bắt buộc phân biệt \code{diem-d} và \code{diem-đ}.

Test Parser Router kiểm tra ma trận routing: PDF searchable sang Docling; PDF scan sang Docling rồi quality gate fail sang MinerU; bảng phức tạp so sánh đầu ra hai parser. Quality gate nhóm A và nhóm B có test riêng; khi structural gate fail, dữ liệu structural cũ bị hủy và chạy lại từ parser thay thế, không trộn kết quả hai parser.

\section{Test temporal và citation}

\begin{sloppypar}
Boundary test chạy tại \code{effective\_from}, \code{effective\_to} trừ một ngày và \code{effective\_to} (upper bound exclusive). Temporal and Amendment Resolver được test với fixture sự kiện AMENDED, PARTIAL\_AMENDED, REPEALED, SUPERSEDED và CORRECTED; mỗi event test kèm một query current và một query historical để khẳng định phiên bản đúng được chọn. Citation verifier có test cho ID ngoài context, provision không tồn tại, document chưa accepted, date invalid, claim không citation, số liệu sai (L4) và claim không được hỗ trợ (L5). Test repair workflow khẳng định mọi nhánh repair tăng cùng counter \code{repair\_attempts}, hết giới hạn kết thúc ở ABSTAIN và không có vòng lặp vô hạn.
\end{sloppypar}

\section{Thiết lập evaluation}

Mỗi run đọc frozen corpus, frozen gold set và variant config; result được align bằng question ID, không dùng \code{zip}; missing result làm run failed. Run và raw artifact bất biến/append-only: mỗi run ghi \code{run\_manifest\_hash}, model IDs, prompt versions, parser versions, token usage, estimated cost và raw results path. Deterministic metric là headline; GPT-5.4 mini snapshot chỉ chấm metric semantic phụ; citation, temporal và numeric metrics được tính bằng code. Dev set dùng để lặp phát triển, validation set dùng để chọn ngưỡng/model/prompt, final test set đóng băng và không dùng để tuning.

\section{Thiết kế thí nghiệm các suite}

\subsection{Suite A - Parser benchmark (P1-P3)}

\pendingresult{Chạy P1 (Docling), P2 (MinerU) và P3 (Parser Router) trên cùng fixture với gold annotation; điền Article/Clause/Point P/R/F1, Short Point Recall, Vietnamese đ) Recall, Parent Context Completeness, Table Preservation, Header/Footer Leakage và Provenance Coverage.}

\subsection{Suite B - Embedding benchmark (E1-E3)}

\pendingresult{Chạy E1 (Gemini Embedding 2), E2 (Jina v5 text-nano) và E3 (Jina v5 text-small) trên development set; điền Recall@10, MRR@10, nDCG@10, latency và cost. Quyết định embedding production dựa trên kết quả này.}

\subsection{Suite C - Retrieval ablation (R1-R10)}

\pendingresult{Chạy các variant R1-R10 trên validation set; điền Recall@5/10/20, MRR@10, nDCG@10 cho từng variant và phân tích tầng nào đóng góp nhiều nhất. Không chọn config bằng final result.}

\subsection{Suite D - Generation và verification ablation (G1-G7)}

\pendingresult{Chạy G1-G7 trên validation set; điền Citation Precision/Recall/F1, Returned Invalid Citation Rate, Numeric Grounding Accuracy, Unsupported Claim Rate, Evidence Completeness và Abstention P/R/F1 cho từng variant.}

\section{Kết quả extraction}

\pendingresult{Điền bảng Suite A: Article/Clause/Point P/R/F1, Short Point Recall, Vietnamese đ) Recall, Parent Context Completeness, Table Preservation, Header/Footer Leakage và Provenance Coverage sau khi review bộ fixture cuối.}

\begin{table}[H]
\caption{Kết quả hierarchy (Suite A) - chờ final evaluation}
\small
\begin{tabularx}{\textwidth}{X c c c c c c c c c}
\toprule
\multirow{2}{*}{\textbf{Phương pháp}} & \multicolumn{3}{c}{\textbf{Article}} & \multicolumn{3}{c}{\textbf{Clause}} & \multicolumn{3}{c}{\textbf{Point}}\\
\cmidrule(lr){2-4}\cmidrule(lr){5-7}\cmidrule(lr){8-10}
 & \textbf{P} & \textbf{R} & \textbf{F1} & \textbf{P} & \textbf{R} & \textbf{F1} & \textbf{P} & \textbf{R} & \textbf{F1}\\
\midrule
P1 Docling & -- & -- & -- & -- & -- & -- & -- & -- & --\\
P2 MinerU & -- & -- & -- & -- & -- & -- & -- & -- & --\\
P3 Parser Router & -- & -- & -- & -- & -- & -- & -- & -- & --\\
\bottomrule
\end{tabularx}
\end{table}

\begin{table}[H]
\caption{Kết quả các chỉ số bổ sung (Suite A) - chờ final evaluation}
\small
\begin{tabularx}{\textwidth}{X c c c c c c}
\toprule
\textbf{Phương pháp} & \textbf{Short Point} & \textbf{VN đ)} & \textbf{Parent Ctx} & \textbf{Table} & \textbf{H/F Leak} & \textbf{Provenance}\\
 & \textbf{Recall} & \textbf{Recall} & \textbf{Completeness} & \textbf{Preservation} & \textbf{(mục tiêu thấp)} & \textbf{Coverage}\\
\midrule
P1 Docling & -- & -- & -- & -- & -- & --\\
P2 MinerU & -- & -- & -- & -- & -- & --\\
P3 Parser Router & -- & -- & -- & -- & -- & --\\
\bottomrule
\end{tabularx}
\end{table}

\section{Kết quả retrieval}

\pendingresult{Chạy retrieval-only trên final test set đã đóng băng. Không chọn config bằng final result.}

\begin{table}[H]
\caption{Kết quả retrieval - chờ final evaluation}
\begin{tabularx}{\textwidth}{c c c c c c}
\toprule
\textbf{Variant} & \textbf{R@5} & \textbf{R@10} & \textbf{MRR@10} & \textbf{nDCG@10} & \textbf{Evidence Set Recall}\\
\midrule
R1 & -- & -- & -- & -- & --\\
R2 & -- & -- & -- & -- & --\\
R3 & -- & -- & -- & -- & --\\
R4 & -- & -- & -- & -- & --\\
R5 & -- & -- & -- & -- & --\\
R6 & -- & -- & -- & -- & --\\
R7 & -- & -- & -- & -- & --\\
R8 & -- & -- & -- & -- & --\\
R9 & -- & -- & -- & -- & --\\
R10 & -- & -- & -- & -- & --\\
\bottomrule
\end{tabularx}
\end{table}

\section{Kết quả temporal và citation}

\pendingresult{Điền kết quả G4-G7, phân biệt Draft Invalid Citation Rate và Returned Invalid Citation Rate; báo Temporal Leakage Rate theo category.}

\begin{table}[H]
\caption{Kết quả temporal và citation - chờ final evaluation}
\begin{tabularx}{\textwidth}{c c c c c c}
\toprule
\textbf{Variant} & \textbf{Temporal Acc.} & \textbf{Cit. P} & \textbf{Cit. R} & \textbf{Cit. F1} & \textbf{Returned Invalid}\\
\midrule
G4 & -- & -- & -- & -- & --\\
G5 & -- & -- & -- & -- & --\\
G6 & -- & -- & -- & -- & --\\
G7 & -- & -- & -- & -- & --\\
\bottomrule
\end{tabularx}
\end{table}

\section{Kết quả generation và abstention}

\begin{table}[H]
\caption{Kết quả generation - chờ final evaluation}
\small
\begin{tabularx}{\textwidth}{c c c c c}
\toprule
\textbf{Variant} & \textbf{Numeric Acc.} & \textbf{Unsupported Claim} & \textbf{Ev. Completeness} & \textbf{Abstention F1}\\
\midrule
G6 & -- & -- & -- & --\\
G7 & -- & -- & -- & --\\
\bottomrule
\end{tabularx}
\end{table}

\section{Latency, tài nguyên và chi phí}

\begin{table}[H]
\caption{Kết quả vận hành - chờ benchmark}
\begin{tabularx}{\textwidth}{X c c c c}
\toprule
\textbf{Cấu hình} & \textbf{Retrieval P95} & \textbf{E2E P95} & \textbf{RAM đỉnh} & \textbf{Cost/query}\\
\midrule
R10 local defense & -- & -- & -- & --\\
\bottomrule
\end{tabularx}
\end{table}

Các ngưỡng mục tiêu của hệ thống: retrieval P95 không quá 2 giây, end-to-end P50 không quá 12 giây, end-to-end P95 không quá 20 giây, search API P95 không quá 3 giây. Đây là mục tiêu kỹ thuật trong môi trường test, không phải kết quả đo.

\section{Phân tích lỗi}

Failure được gán một root cause trong các nhóm parser routing, parser OCR/layout, IR normalization, hierarchy extraction, provenance, dense retrieval, sparse retrieval, fusion, query analysis, query expansion, reranking, context expansion, evidence gate, temporal filter, generation, citation selection, numeric grounding, claim support, abstention, gold set và provider.

\pendingresult{Sau final run, chọn 10-15 lỗi đại diện, mô tả triệu chứng, nguyên nhân, ảnh hưởng và hướng sửa. Không loại câu khó khỏi báo cáo.}

\section{Diễn giải kết quả}

Khi có số liệu, phần diễn giải phải trả lời trực tiếp RQ1-RQ6. Cần báo cả absolute difference, per-category difference và failure trade-off. Ví dụ, verifier sáu tầng có thể giảm invalid citation nhưng tăng over-abstention; đây là trade-off cần báo cáo thay vì chỉ trình bày một điểm tổng hợp. Mọi báo cáo aggregate bắt buộc kèm phân rã theo category (CURRENT, HISTORICAL, COMPARISON, MULTI\_PROVISION, ...) vì một điểm trung bình có thể che lỗi của nhóm câu quan trọng.

\section{Tiêu chí release}

Release candidate chỉ được chấp nhận khi migration từ database rỗng, restore từ release bundle, ba luồng E2E (hiện hành, lịch sử, so sánh), abstention flow, contract Returned Invalid Citation Rate = 0 và regression subset đều đạt. Báo cáo và slide phải dùng cùng evaluation artifact để tránh số liệu không đồng nhất. Trước khi nộp, chuyển \code{\textbackslash draftmodefalse} và điền mọi bảng kết quả từ evaluation run đã đóng băng.

% ============================================================
\chapter{Kết Luận và Hướng Phát Triển}
% ============================================================

\section{Kết luận}

Khóa luận đề xuất một kiến trúc RAG chuyên biệt cho pháp luật giao thông đường bộ Việt Nam, tập trung vào structure-aware parsing, mô hình quan hệ tham chiếu, temporal retrieval và verifiable citations. Parser Router và Canonical Document IR cô lập tầng parsing khỏi phân tích pháp lý; Legal Structure Extractor xử lý phân cấp pháp luật Việt Nam với nhãn Điểm tiếng Việt; PostgreSQL quản lý version, quan hệ và review; Qdrant thực hiện retrieval đa tầng với exact lookup, dense, sparse và RRF; LangGraph điều phối workflow có giới hạn; generator chỉ được chọn \code{provision\_id} từ context; verifier sáu tầng quyết định trả verified answer hay abstain.

Thiết kế này loại bỏ các thành phần khó kiểm soát của phương án cũ như web fallback, citation tự do, index retrieval tách rời và vòng lặp điều phối/sửa lỗi không giới hạn. Hệ thống được tổ chức để một lỗi có thể được truy ngược về parser, extraction, retrieval, temporal metadata, generation hoặc verification. Langfuse cung cấp observability nhưng không nằm trên đường tới hạn tính đúng đắn; RAGFlow chỉ đóng vai trò baseline so sánh bên ngoài.

\pendingresult{Trước khi nộp, bổ sung kết luận định lượng dựa trên kết quả bốn suite A-D: phương pháp nào tốt nhất, mức cải thiện, returned invalid citation, temporal accuracy, evidence completeness, latency và chi phí.}

\section{Ưu điểm}

\begin{itemize}
  \item mô hình hóa hierarchy, quan hệ tham chiếu và temporal interval như dữ liệu nghiệp vụ trong PostgreSQL;
  \item Parser Router và Canonical Document IR cho phép đánh giá và thay parser mà không viết lại Legal Structure Extractor;
  \item nhận diện nhãn Điểm tiếng Việt \code{d)} và \code{đ)} với stable ID tái tạo được;
  \item provenance đến page/bounding box và truy vết về IR;
  \item retrieval đa tầng trong một vector database với exact lookup, dense, sparse, RRF và reranker;
  \item evidence planning và Evidence Completeness Gate chặn câu trả lời nửa vời;
  \item citation-by-ID và verification sáu tầng với invariant Returned Invalid Citation Rate = 0;
  \item verified-or-abstain contract với failure-aware repair có giới hạn;
  \item evaluation có split, hash và raw artifacts bất biến;
  \item local-first deployment với Docker Compose bảy service và restore plan.
\end{itemize}

\section{Hạn chế}

\begin{itemize}
  \item corpus nhỏ và chỉ tập trung giao thông đường bộ;
  \item partial amendment mapping cần nhiều review thủ công;
  \item một reviewer chính có thể tạo annotation bias;
  \item model API và stable alias vẫn có thể drift giữa các lần chạy;
  \item claim support judge chưa phải legal expert;
  \item MinerU VLM/hybrid backend không khả thi trên máy cá nhân, chỉ dùng pipeline CPU;
  \item hệ thống là công cụ hỗ trợ tra cứu, không thay tư vấn pháp lý.
\end{itemize}

\section{Hướng phát triển}

Hướng ngắn hạn gồm admin review UI, conversation history, diff viewer, calibrated abstention và khóa ngưỡng reranker/evidence gate sau khi có kết quả Suite C-D. Hướng trung hạn gồm provision-level amendment engine, domain expert workflow, mở rộng baseline RAGFlow sang các dataset mới và multilingual query. Hướng dài hạn gồm citation entailment model tiếng Việt, mở rộng sang các domain pháp luật khác, nghiên cứu late-interaction/ColBERT-style reranking và tái sử dụng gold set cho các bộ dữ liệu công khai.

\section{Bài học phương pháp}

Các nhận định như parser nào tốt hơn, hybrid retrieval tốt hơn dense-only, reranker cải thiện chất lượng hay verifier giảm invalid citation chỉ được kết luận sau thực nghiệm trên gold set đã review, không dựa trên tuyên bố nhà cung cấp. Bài học cốt lõi của thiết kế là chất lượng hệ thống pháp luật phụ thuộc đồng thời vào dữ liệu có nguồn được review, ranh giới pháp lý trùng ranh giới trích dẫn, retrieval áp đúng phiên bản hiệu lực, quan hệ tham chiếu được mở rộng có giới hạn và contract không cho phép answer chưa verify đi qua.

% ============================================================
% REFERENCES
% ============================================================
\clearpage
\printbibliography[heading=bibintoc,title={Tài Liệu Tham Khảo}]

% ============================================================
% APPENDICES
% ============================================================
\appendix

\chapter{Hướng Dẫn Chạy Demo}

\section{Khởi động hạ tầng}

Theo quy định trình bày khóa luận của Khoa Kỹ thuật - Công nghệ, Trường Đại học Tây Đô \cite{taydo2024format}, mọi lệnh \code{docker compose} chạy từ repo root và kèm \code{--project-directory .}.

\begin{lstlisting}[caption={Khởi động defense mode}]
cp deploy/env/release.env.example .env
# Dien API key va secret

docker compose \
  --project-directory . -f deploy/compose/compose.release.yml \
  up -d

docker compose --project-directory . \
  -f deploy/compose/compose.release.yml wait backend frontend

curl --fail http://localhost:8000/api/v1/health/ready
\end{lstlisting}

\section{Kịch bản demo}

\begin{enumerate}
  \item Hỏi quy định hiện hành.
  \item Mở citation và source passage kèm trang.
  \item Hỏi cùng hành vi tại thời điểm lịch sử.
  \item So sánh trước và sau một mốc hiệu lực.
  \item Hỏi ngoài corpus để minh họa abstention.
  \item Hỏi câu đa bằng chứng (mức phạt + điểm trừ) để minh họa evidence completeness.
  \item Mở bảng kết quả evaluation bốn suite A-D.
\end{enumerate}

\chapter{Schema Gold Set}

\begin{lstlisting}[language=json,caption={Gold record schema rút gọn}]
{
  "id": "Q001",
  "split": "FINAL_TEST",
  "question": "Năm 2023 xe máy vượt đèn đỏ bị xử lý thế nào?",
  "category": "HISTORICAL",
  "query_date": "2023-07-01",
  "vehicle_type": "MOTORCYCLE",
  "expected_provision_ids": ["..."],
  "acceptable_provision_ids": ["..."],
  "required_evidence": ["monetary_penalty"],
  "must_include_facts": ["..."],
  "must_not_include_facts": ["..."],
  "temporal_metadata": {
    "query_date": "2023-07-01",
    "date_source": "canonical_date"
  },
  "review_status": "REVIEWED",
  "reviewed_by": "reviewer-01",
  "gold_version": "gold-v1",
  "hash": "sha256:..."
}
\end{lstlisting}

Mọi định danh pháp lý trong ví dụ là placeholder tổng hợp; ID thật được xác định theo quy trình tạo gold và validate trong PostgreSQL. Câu lịch sử năm 2023 phải trỏ tới provision của văn bản hợp lệ tại mốc 2023, không trỏ tới văn bản ban hành sau mốc hỏi.

\chapter{Checklist Trước Khi Nộp}

\begin{itemize}
  \item[$\square$] Điền tên giảng viên hướng dẫn và mã số sinh viên.
  \item[$\square$] Chèn logo chính thức và screenshot giao diện.
  \item[$\square$] Chuyển \code{\textbackslash draftmodefalse} sau khi có final results.
  \item[$\square$] Điền các bảng kết quả Chương 6 từ evaluation artifact.
  \item[$\square$] Cập nhật Tóm tắt, Abstract và Kết luận bằng số liệu thật.
  \item[$\square$] Kiểm tra mọi citation trong bibliography.
  \item[$\square$] Kiểm tra danh mục hình, bảng và mục lục.
  \item[$\square$] Build bằng XeLaTeX + Biber từ clean workspace.
  \item[$\square$] Render PDF và kiểm tra clipping, font, table overflow.
  \item[$\square$] Đồng bộ số liệu giữa báo cáo và slide.
  \item[$\square$] Tạo release bundle và video backup.
\end{itemize}

\end{document}
